%%% The main file. It contains definitions of basic parameters and includes all other parts.

%%% Settings for single-side (simplex) printing
%% Margins: left 40mm, right 25mm, top and bottom 25mm
%% (but beware, LaTeX adds 1in implicitly)
%\documentclass[12pt,a4paper,usenames]{report}
%\setlength\textwidth{145mm}
%\setlength\textheight{247mm}
%\setlength\oddsidemargin{15mm}
%\setlength\evensidemargin{15mm}
%\setlength\topmargin{0mm}
%\setlength\headsep{0mm}
%\setlength\headheight{0mm}
%% \openright makes the following text appear on a right-hand page
%\let\openright=\clearpage

%% Settings for two-sided (duplex) printing
 \documentclass[12pt,a4paper,twoside,openright]{report}
 \setlength\textwidth{145mm}
 \setlength\textheight{247mm}
 \setlength\oddsidemargin{14.2mm}
 \setlength\evensidemargin{0mm}
 \setlength\topmargin{0mm}
 \setlength\headsep{0mm}
 \setlength\headheight{0mm}
 \let\openright=\cleardoublepage

 \raggedbottom

%%disable overfull lines and disable hphenation
%%\usepackage[none]{hyphenat}
%\usepackage[english]{babel}
%\selectlanguage{english}
%\usepackage{hyphenat}
%\allowhyphens
%%\overfullrule=0mm
\hfuzz=5mm
\setlength{\emergencystretch}{10pt}


\interfootnotelinepenalty=10000


\DeclareTextFontCommand{\ttt}{\ttfamily\hyphenchar\font=45\relax}
\DeclareTextFontCommand{\trm}{\rmfamily\hyphenchar\font=45\relax}

%\renewcommand{\baselinestretch}{1.5}

%% Character encoding: usually latin2, cp1250 or utf8:
\usepackage[utf8]{inputenc}


%% Further useful packages (included in most LaTeX distributions)
%\usepackage{amsmath}        % extensions for typesetting of math
\usepackage{amsfonts}       % math fonts
\usepackage{amsthm}         % theorems, definitions, etc.
\usepackage{thmtools}
\usepackage{bbding}         % various symbols (squares, asterisks, scissors, ...)
\usepackage{bm}             % boldface symbols (\bm)
\usepackage{graphicx}       % embedding of pictures
\usepackage{fancyvrb}       % improved verbatim environment
\usepackage{natbib}         % citation style AUTHOR (YEAR), or AUTHOR [NUMBER]
\usepackage[nottoc]{tocbibind} % makes sure that bibliography and the lists
			    % of figures/tables are included in the table
			    % of contents
\usepackage{dcolumn}        % improved alignment of table columns
\usepackage{booktabs}       % improved horizontal lines in tables
\usepackage{paralist}       % improved enumerate and itemize
\usepackage[usenames]{xcolor}  % typesetting in color

\usepackage{tikz}
\usepackage{placeins}
%\usepackage{animate}
\usepackage{float}
\usetikzlibrary{shapes,snakes,arrows,chains}
\usepackage{amsmath}
\usepackage{rotating}

\usepackage{dirtytalk}

\usepackage{dirtree}
\usepackage{url}

\usepackage{longtable}
%%% Basic information on the thesis

% Thesis title in English (exactly as in the formal assignment)
\def\ThesisTitle{SIMD code generator}

% Author of the thesis
\def\ThesisAuthor{Karel Tuček}

% Year when the thesis is submitted
\def\YearSubmitted{2016}

% Name of the department or institute, where the work was officially assigned
% (according to the Organizational Structure of MFF UK in English,
% or a full name of a department outside MFF)
\def\Department{Department of Software Engineering}

% Is it a department (katedra), or an institute (ústav)?
\def\DeptType{Department}

% Thesis supervisor: name, surname and titles
\def\Supervisor{RNDr. David Bednárek, Ph.D.}

% Supervisor's department (again according to Organizational structure of MFF)
\def\SupervisorsDepartment{Department of Software Engineering}

% Study programme and specialization
\def\StudyProgramme{Computer Science}
\def\StudyBranch{General Computer Science}

% An optional dedication: you can thank whomever you wish (your supervisor,
% consultant, a person who lent the software, etc.)
\def\Dedication{%
}

% Abstract (recommended length around 80-200 words; this is not a copy of your thesis assignment!)
\def\Abstract{%
The center of our interest is a problem of pipelined realisation of a special case of data processing networks. These realisations are supposed to realise some computations on series of independent data sets while utilizing SIMD instructions. The aim of this paper is to theoretically investigate the possibilities and the problems of employment of control flow in these networks and also to implement a general framework suitable for generation of these realisations.

The main idea is utilisation of an algorithm crawling over partitions of a network factorised with respect to its control flow. Our idea is that SIMD parallelism should take place on the same instruction realised across multiple data sets.

We illustrate the problems relevant to employment of branching and loops in these networks. We especially discuss a problem of data ordering and also provide relevant proofs. In the analytical part, we show implementation of a general framework which we believe to be suitable for processing of these networks. We also provide examples utilising Intel's SIMD Streaming Extensions.
}

% 3 to 5 keywords (recommended), each enclosed in curly braces
\def\Keywords{%
  {Processing networks} {SIMD} {Parallelism}
}

%% The hyperref package for clickable links in PDF and also for storing
%% metadata to PDF (including the table of contents).
\usepackage[pdftex,unicode]{hyperref}   % Must follow all other packages
\hypersetup{breaklinks=true}
\hypersetup{pdftitle={\ThesisTitle}}
\hypersetup{pdfauthor={\ThesisAuthor}}
\hypersetup{pdfkeywords=\Keywords}
\hypersetup{urlcolor=blue}

% Definitions of macros (see description inside)
%%% This file contains definitions of various useful macros and environments %%%
%%% Please add more macros here instead of cluttering other files with them. %%%

%%% Minor tweaks of style

% These macros employ a little dirty trick to convince LaTeX to typeset
% chapter headings sanely, without lots of empty space above them.
% Feel free to ignore.
\makeatletter
\def\@makechapterhead#1{
  {\parindent \z@ \raggedright \normalfont
   \Huge\bfseries \thechapter. #1
   \par\nobreak
   \vskip 20\p@
}}
\def\@makeschapterhead#1{
  {\parindent \z@ \raggedright \normalfont
   \Huge\bfseries #1
   \par\nobreak
   \vskip 20\p@
}}
\makeatother

% This macro defines a chapter, which is not numbered, but is included
% in the table of contents.
\def\chapwithtoc#1{
\chapter*{#1}
\addcontentsline{toc}{chapter}{#1}
}

% Draw black "slugs" whenever a line overflows, so that we can spot it easily.
\overfullrule=1mm

%%% Macros for definitions, theorems, claims, examples, ... (requires amsthm package)

\theoremstyle{plain}
\newtheorem{thm}{Theorem}
\newtheorem{theorem}{Theorem}
\newtheorem{lemma}[thm]{Lemma}
\newtheorem{claim}[thm]{Claim}

\newtheorem{observation}{Observation}
\newtheorem{problem}{Problem}
\newtheorem{define}{Definition}
\newtheorem{alg}{Algorithm}
\newtheorem*{defstyle}{}

\theoremstyle{remark}
\newtheorem*{rem}{Remark}
\newtheorem*{cor}{Corollary}
\newtheorem*{example}{Example}

\newcommand{\mybeginfigloose}{\begin{figure}[h]}
\newcommand{\mybeginfig}{\begin{figure}[H]}
\newcommand{\myendfig}[2]{\caption{#2}\label{fig:#1}\end{figure}}

\newcommand{\mybegindef}[2]{\begin{define}[#2]\label{def:#1}}
\newcommand{\myenddef}{\end{define}}

\newcommand{\mybeginalg}[2]{\begin{alg}[#2]\label{alg:#1}}
\newcommand{\myendalg}{\end{alg}}

\newcommand{\mybeginprob}[2]{\begin{problem}[#2]\label{pro:#1}}
\newcommand{\myendprob}{\end{problem}}

\newcommand{\mybeginobs}[2]{\begin{observation}[#2]\label{obs:#1}}
\newcommand{\myendobs}{\end{observation}}

\newcommand{\mybeginclaim}[2]{\begin{claim}[#2]\label{cla:#1}}
\newcommand{\myendclaim}{\end{claim}}

%\newcommand{\myexpr}[1]{ $\begin{gathered}[t]\begin{aligned} #1 \end{aligned}\end{gathered}$ }

\newcommand{\myquote}[2]{\\ \\ \centerline{\sl{``#1''}\cite{#2}}}
\newcommand{\myloosequote}[2]{{\sl{``#1''}\cite{#2}}}

\newcommand{\algref}[1]{\begin{alg}[#1]\end{alg}}

\newcommand{\myexpr}[1]{ $ #1 $ }

\newcommand{\allowpg}[1]{ \end{code}\ \begin{code} }

\newcommand{\gt}{ > }
\newcommand{\lt}{ < }


%tree items
%\usetikzlibrary{shadows}
%\newcommand{\treelist}{\tikz[overlay]\draw(-.2,-.2)--(-.2,.5) (-.2,.15)--(.1,.15);}
%\newcommand{\treeitem}{\item[\treelist]}
%\newcommand{\treeendlist}{\tikz[overlay]\draw(-.2,-.2)--(-.2,.15) --(.1,.15);}
%\newcommand{\treeend}{\item[\treeendlist]}

%\newcommand{\begincodefig}[2]{\begin{figure}[H]\caption{#2}\label{fig:#1}\begin{code}}
%\newcommand{\endcodefig}{\end{code}\end{figure}}


%moje pridavky
\newcommand{\graphplain}[1]{\begin{figure}[H]\centering\input{contents/graphs/final/#1}\end{figure}}
\newcommand{\graph}[2]{\begin{figure}[H]\centering\input{contents/graphs/final/#1}\caption{#2}\label{fig:#1}\end{figure}}
%\newcommand{\graph}[1]{\begin{figure}[H]\centering\resizebox{.9\linewidth}{!}{\input{contents/graphs/final/#1}}\end{figure}}
  \newcommand{\graphn}[1]{\begin{sidewaysfigure}\resizebox{1\linewidth}{!}{\input{contents/graphs/final/#1}}\end{sidewaysfigure}}
%k\newcommand{\graph}[1]{\input{contents/graphs/final/#1}\end{figure}}
%k\resizebox{.9\linewidth}{!}{\input{plot.tex}}
\newcommand{\parspace}{\vspace{1\baselineskip}}

%%% An environment for proofs

%%% FIXME %%% \newenvironment{proof}{
%%% FIXME %%%   \par\medskip\noindent
%%% FIXME %%%   \textit{Proof}.
%%% FIXME %%% }{
%%% FIXME %%% \newline
%%% FIXME %%% \rightline{$\square$}  % or \SquareCastShadowBottomRight from bbding package
%%% FIXME %%% }

%%% An environment for typesetting of program code and input/output
%%% of programs. (Requires the fancyvrb package -- fancy verbatim.)


\DefineVerbatimEnvironment{loosecode}{Verbatim}{fontsize=\small, frame=leftline, fontshape=n, samepage=false} 
\DefineVerbatimEnvironment{code}{Verbatim}{fontsize=\small, frame=leftline, fontshape=n, samepage=true} 
\newcommand{\maybreak}[1]{\allowbreak}

\newcommand{\incode}[1]{ \\ \centerline{\texttt{#1}} \\ }

%%% The field of all real and natural numbers
\newcommand{\R}{\mathbb{R}}
\newcommand{\N}{\mathbb{N}}
\newcommand{\arrow}{\rightarrow}
\newcommand{\powerset}{\raisebox{.15\baselineskip}{\Large\ensuremath{\wp}}}

%%% Useful operators for statistics and probability
\DeclareMathOperator{\pr}{\textsf{P}}
\DeclareMathOperator{\E}{\textsf{E}\,}
\DeclareMathOperator{\var}{\textrm{var}}
\DeclareMathOperator{\sd}{\textrm{sd}}

%%% Transposition of a vector/matrix
\newcommand{\T}[1]{#1^\top}

%%% Various math goodies
\newcommand{\goto}{\rightarrow}
\newcommand{\gotop}{\stackrel{P}{\longrightarrow}}
\newcommand{\maon}[1]{o(n^{#1})}
\newcommand{\abs}[1]{\left|{#1}\right|}
\newcommand{\dint}{\int_0^\tau\!\!\int_0^\tau}
\newcommand{\isqr}[1]{\frac{1}{\sqrt{#1}}}

\newcommand{\uv}[1]{``#1''}

%%% Various table goodies
\newcommand{\pulrad}[1]{\raisebox{1.5ex}[0pt]{#1}}
\newcommand{\mc}[1]{\multicolumn{1}{c}{#1}}


% Title page and various mandatory informational pages
\begin{document}
%%% Title page of the thesis and other mandatory pages

%%% Title page of the thesis

\pagestyle{empty}
\hypersetup{pageanchor=false}
\begin{center}

\large

Charles University in Prague

\medskip

Faculty of Mathematics and Physics

\vfill

{\bf\Large BACHELOR THESIS}

\vfill

\centerline{\mbox{\includegraphics[width=60mm]{./img/logo.pdf}}}

\vfill
\vspace{5mm}

{\LARGE\ThesisAuthor}

\vspace{15mm}

{\LARGE\bfseries\ThesisTitle}

\vfill

\Department

\vfill

\begin{tabular}{rl}

Supervisor of the bachelor thesis: & \Supervisor \\
\noalign{\vspace{2mm}}
Study programme: & \StudyProgramme \\
\noalign{\vspace{2mm}}
Study branch: & \StudyBranch \\
\end{tabular}

\vfill

% Zde doplňte rok
Prague \YearSubmitted

\end{center}

\newpage

%%% Here should be a bound sheet included -- a signed copy of the "bachelor
%%% thesis assignment". This assignment is NOT a part of the electronic
%%% version of the thesis. DO NOT SCAN.

%%% A page with a solemn declaration to the bachelor thesis

\openright
\hypersetup{pageanchor=true}
\pagestyle{plain}
\pagenumbering{roman}
\vglue 0pt plus 1fill

\noindent
I declare that I carried out this bachelor thesis independently, and only with the cited
sources, literature and other professional sources.

\medskip\noindent
I understand that my work relates to the rights and obligations under the Act No.~121/2000 Sb.,
the Copyright Act, as amended, in particular the fact that the Charles
University in Prague has the right to conclude a license agreement on the use of this
work as a school work pursuant to Section 60 subsection 1 of the Copyright Act.

\vspace{10mm}

\hbox{\hbox to 0.5\hsize{%
In ........ date ............	% FIXME!
\hss}\hbox to 0.5\hsize{%
signature of the author
\hss}}

\vspace{20mm}
\newpage

%%% Mandatory information page of the thesis

\openright

\vbox to 0.5\vsize{
\setlength\parindent{0mm}
\setlength\parskip{5mm}

Title:
\ThesisTitle

Author:
\ThesisAuthor

\DeptType:
\Department

Supervisor:
\Supervisor, \SupervisorsDepartment

Abstract:
\Abstract

Keywords:
\Keywords

\vss}

\newpage

%%% Dedication

\openright

\noindent
\Dedication

\newpage

\openright
\pagestyle{plain}
\pagenumbering{arabic}
\setcounter{page}{1}


%%% A page with automatically generated table of contents of the bachelor thesis

\tableofcontents

%%% Each chapter is kept in a separate file
\setlength{\parskip}{4mm plus2mm minus2mm}
\chapter*{Introduction}
\addcontentsline{toc}{chapter}{Introduction}


\chapter{Introduction}
  %\chapter*{Introduction}
%\addcontentsline{toc}{chapter}{Introduction}

\section{Overview}

TODO: some broader motivation (What?)

The main goal of this work is implementation of a framework which may serve as the backmost part of a (trans) compiler\footnote{Trans-compiler is a compiler which translates source code into some other programmememing language or another intermediate form.} and which would provide some specific parallelism-related features.\footnote{Do not panic!} We also wish to investigate the problem of employment of control flow\footnote{The term \emph{control flow} is generally used for all mechanisms that use some form of the goto instruction, like branching, loops, switches (\dots). Although the general meaning encompasses function calls, we will use this term only for branching and loops.} in parallelized environments of a specific form.

\section{Problem introduction, input examples}

Our problem is basically turning an intermediate code represented by some oriented graph and stripped of some type of dependencies into a fragment of code in some target language while utilising some intrinsic\footnote{Providing direct access to instructions of the target platform} SIMD\footnote{\emph{Single-Instruction Multiple-Data} instructions are instructions which take multiple independent sets of input values and perform an operation parallelly on all of them. E.g., performing an element-wise addition of two vectors.} extension. The aim is to achieve a situation in which multiple evaluations of the same computation are processed at the same time employing the mentioned fine-grained-parallelism. 


We assume input in a form of a graph annotated by some operations. We assume this input to be either written by hand or produced by a specialised software similar to a compiler back-end. We expect this software to identify, extract and preprocess code which is suitable for fine-grained parallelism.  An example of this software may be Parallax \cite{parallax}. We do not provide any integration at the moment since Parallax is still under development and exact interfaces are still unclear.


Since textual description of the problem may be too unclear, we provide some examples.

\graph{adder0}{Simple adder computation}

This graph represents a computation which takes values from two (external) streams, sums them and produces another (external) stream. The LD operations represent some data input for our algorithm, i.e., the LD may evaluate to some effectful\footnote{When saying \emph{effect}, we mean (\emph{impure}) side effects of operations or concealed data flow. We use this word the same way functional programmememers do.} expressions. The ADD operation represents addition. The ST operation represents an effectfull expression which hands the produced values over to a consumer. 


So, this graph surely could be understood as a representation of the C function given in figure \ref{fig:adder0code}. However, semantical meaning we wish to assign to these graphs is a bit different. A more accurate interpretation understands this graph to be a computation\footnote{We will try to use this term consistently with its standard use in functional programmememing, i.e., as a data channel which transforms the data flowing through. However, at some places we may not manage to keep this promise since in some cases we lack more exact terms.} which repeatedly invokes the function, with possible overlaps in time. In other words, we have performed the \emph{lift} transformation, which is well known in functional programmememing.

\mybeginfig
\begin{code}
int f(int a, int b)
{
  return a+b;
}
\end{code}
\myendfig{adder0code}{C fragment corresponding to figure \ref{fig:adder0}}

In this explanation we have used the word \emph{stream} strictly for external streams (managed by some third party). Every external stream stands merely for an input or an output for our computation. This makes sense since every computation, represented by a graph, may be seen as a building block of some higher level pipeline. Such pipelines may be (and are supposed to be) constructed by the specialised software mentioned above. Another thing is that the internal functionality of these computations also resembles (or will resemble) pipelines, and therefore we will use the word \emph{pipeline} for these computations and the word \emph{stream} for edges in our graphs. We will also talk about these edges as of \emph{queues}. 


From this point of view, the load and store instructions may be perceived as operations which merely transform external streams into internal streams. Note that although we needed to make this clear, we will not talk about or work with external streams anymore. Loads and stores will be just data sources and sinks for us. Also, we will use dashed lines no more.


Output of our framework is supposed to be a code fragment which receives some input data in form somehow equivalent to a table such that:
\begin{itemize}
  \item Columns of this table correspond to input operations in the graph. Any column also may be seen as a single data stream or as a set of synchronous data streams. 
  \item Every line describes input data for one calculation. (For instance, input data for one invocation of a function corresponding to a graph)
  \item Different rows represent independent calculations (but still defined by the same graph). 
\end{itemize}
The code then produces some results per every row of the table, effectively adding some new columns.

  \mybeginfig
\begin{center}
  \begin{tabular}{c|c||c}
    a (load)&b (load)&d (store)\\
  \hline
  1&2&3\\
  4&5&9\\
  7&8&15\\
  3&2&5\\
  \end{tabular}
\end{center}
\myendfig{adder0data}{Example data input and output for example figure \ref{fig:adder0} }

The produced code fragment (e.g., C code) may, for instance, take the columns $a$ and $b$ and produce the column $d$. 

%, since full implementation exceeds the scope of this work,

The second problem, which we wish to touch on a mostly theoretical level, is employment of control flow in these networks. Consider the following C function. 

\mybeginfig
\begin{code}
int f(int a)
{
  if(a % 2 == 1)
    a = a*7;
  else
    while ( a < 100 )
      a = a+3;
  return a;
}
\end{code}
\myendfig{introcfcode}{Example C fragment employing control flow}

Control flow may be represented in different ways. Since we wish to handle data in a streamlined fashion, splitting and merging data streams proposes itself. We will restrict our attention to control flow in form of branching and loops represented in this fashion. To be more precise, we restrict ourselves to investigation of control flow which actually splits data into multiple streams. We also assume nesting of the control flow constructs.

\begin{rem}
  What we will not develop is the \emph{`compute all branches on all data and select the desired results'}\footnote{We believe that every serious reader of this work knows what we mean. We provide a bit more detailed explanation in section \ref{sec:branching}, but a casual reader should feel free to skip this reference safely.} version of the problem. One of the reasons is that it has been satisfactorily investigated before (e.g., \cite{who_knows}). Another, more pragmatic, reason is that our framework is capable of employing this mechanism in means of standard arithmetic-like operations as long as the input is properly preprocessed.
\end{rem}


The control flow thus will be represented by some form of split and merge operations. These may be understood as operations splitting (and merging) data streams in a non-1:1 ratio. In addition, we may use (and require in our input) some special condition operations. For instance, the provided C function may be represented by the graph in figure \ref{fig:introcf}.

\graph{introcf}{Graph representation of figure \ref{fig:introcfcode}}

\begin{description}
  \item[ST,LD] - again some effectfull IO expressions.
  \item[LT] - \emph{less than} condition with 100 as its second operand.
  \item[ADD,MOD,MUL] - addition, modulo and multiplication, all used with constant second operands.
  \item[SPLIT] - a special control flow operation which takes a value and a boolean and sends the value either to its right or left output depending on the supplemented boolean.
  \item[MERGE] - exact opposite of the SPLIT operation.
  \item[CONDITION] - a special control flow operation which manages loops. 
\end{description}

  While we believe we have sufficiently illustrated problems of our interest, we have been very vague about details of context in which our framework is assumed to work (according to our assignment). To put this right, we provide the following diagram.

\graph{integration}{Diagram describing assumed integration of our framework.}

\begin{description}
  \item[Specialised software] such as Parallax takes representation of a programmeme and transforms it into some graph representation. Furthermore, it cuts the graph into small chunks and extracts parallelizable excerpts, which it sends to our framework. Besides, it creates some representation of a higher level pipeline and some control C code.
  \item[Graphs] are the graphs described in the previous subsection. We assume these graphs to be in form of xml-formatted files or streams.
  \item[Description of instructions] is some definition of the instruction set of the target language and SIMD extension. 
  \item[Our framework] receives these small parallelizable graphs and transforms them into C modules which may serve as building blocks for higher-level pipelines. These C modules have stream interfaces.
  \item[A compiler] is used to compile all sources requiring compilation.
  \item[Executables] are either more or less conventional executable files which contain all introduced computations.
  \item[Control graph] is some description of a task for the distributed environment.
  \item[Distributed environment] such as Bobox \cite{bobox} is an environment which finally composes all computations and other auxiliary C modules into a final pipeline and runs it. 
\end{description}

\section{Overview of employed methods}

In this paper, we begin by showing that pipelines without control flow may be easily transformed into the anticipated form. We produce these results by generating vertices of the input graph in topological order. This way, the output of every vertex gets stored in a unique variable. 

  
Next thing we show is that every such pipeline may be realised by SIMD instructions only (as long as there is sufficiently complete SIMD set available). The first step is fusing corresponding operations across multiple invocations of the same computation into SIMD instructions. The second step is adding some data conversions on edges which connect instructions of different vector width or alignment. This way we can process data in packs of $LCM$\footnote{Least Common Multiple} of widths of the used instructions. This way instructions may be fused into blocks which may, for instance, somehow correspond to the figure \ref{fig:introfusion}:

\mybeginfig
\begin{center}
  \begin{tabular}{c|c|c|c|c|c|c|c|c|}
    \cline{2-9}
    & \multicolumn{8}{c|}{index of processed row}\\
    \cline{2-9}
    some operation in G & 1 & 2 & 3 & 4 & 5 & 6 & 7 & 8 \\
    \cline{2-9}
    another operation in G & \multicolumn{2}{c|}{1-2} & \multicolumn{2}{c|}{3-4} & \multicolumn{2}{c|}{5-6} & \multicolumn{2}{c|}{7-8} \\
    \cline{2-9}
    \dots & \multicolumn{4}{c|}{1-4} & \multicolumn{4}{c|}{5-8}\\
    \cline{2-9}
    \dots & \multicolumn{2}{c|}{1-2} & \multicolumn{2}{c|}{3-4} & \multicolumn{2}{c|}{5-6} & \multicolumn{2}{c|}{7-8}\\
    \cline{2-9}
    \dots & \multicolumn{8}{c|}{1-8} \\
    \cline{2-9}
    \dots & 1 & 2 & 3 & 4 & 5 & 6 & 7 & 8\\
    \cline{2-9}
  \end{tabular}
\end{center}
\myendfig{introfusion}{Example illustrating how operations may ba joined across different invocations of a computation.}


  Regarding control flow, our approach is to cut the input graph into partitions (basic blocks)\footnote{The term \emph{basic-block} is often used in theory of compilers, meaning consistent sequences of instructions without any control flow. Thus, any procedure may be told to consist of basic blocks and goto instructions.} which process data regularly. That means cutting all control flow nodes into multiple parts and introducing some new data flow between these by some sort of effects. This way we obtain a new graph which consists of basic blocks and control flow edges. We further modify this result to obtain an acyclic graph (or more precisely a graph with some sort of topological ordering\footnote{Ordering of an acyclic graph, which satisfies dependencies represented by edges. This ordering typically represents an order in which steps of some process may be carried out, e.g., steps in the process of cooking a meal. See \cite{kapitoly} for exact definition.}). 

\graph{introfactor}{Graph from figure \ref{fig:introcf} split into partitions suitable for processal}

  After obtaining an acyclic graph consisting of basic blocks, we try to find an algorithm which crawls this new graph and gradually pushes data through the basic blocks. We use (constant-size) buffers \footnote{Memory area which is used to compensate for differences in the rate of flow of data between components of a computer system (\cite{thesarus}). We use the terms \emph{queue} and \emph{buffer} almost interchangeably.} to store intermediate results between these basic blocks.

  %http://www.thefreedictionary.com/buffer

\begin{rem}
  To be precise we allow cycles in a specific form which does not void the \emph{semantical} meaning of topological ordering. This may sound a bit strange, but it is exactly what happens when we use loops in programmememing languages.
\end{rem}

  We should also make clear that our purpose is not to study low-level problematics of compiler construction (these are also well-understood \cite{compiler_theory}). Although we wish to have all intermediate results allocated into CPU registers, we do not deal with register allocation nor we do restrict our register resources. In many aspects, we heavily rely on standard features of target compilers, such as on register allocation or dead code elimination.

To address the question of assumptions for analytical part of this work, we should state the following points:
\begin{itemize}
  \item We assume to have at least $O(n)$ space available in CPU registers with respect to size of input graphs. The actual amount of space needed and available is left up to the provider of our input and to the target compiler. Thus, it is responsibility of the provider to carve up the input into pieces which the target compiler can fit into registers (or into the first level cache) if the provider wishes them to be fitted there.
  \item We wish to minimise the amount of transfers between RAM and CPU.
  \item We wish to utilise fast aligned stores and loads.
  \item Maximal utilisation of the SIMD extension is preferred. Even if the cost needed for data conversions exceeds actual performance gains.
  \item These are part of our assignment.
\end{itemize}

\subsubsection{Summary of our results}

  As we have already stated, we show how graphs without control flow may be solved while utilising a SIMD extension on all data except for some epilogue (e.g. the last few rows of data which do not suffice to fill entire vector). 


  Regarding control flow, there is a problem of data ordering (i.e., requirement that all data should be processed and outputted in the original input order). This problem is discussed in chapters dedicated to control flow .  Furthermore, we discuss the problem of loops, which turns out to be more difficult than originally thought.  


We show two different approaches to this problem - one which preserves the order of the data and one which does not. 


  The first proposed solution, which preserves ordering, may possibly use SIMD instructions on all data in a pipeline without loops. This solution, unfortunately, relies on pull semantics, which allows processing of basic blocks without parallelism. Hence, we have almost no guarantee that SIMD instructions will be employed in runtime. Also, this solution does not work well with loops due to the ordering requirement. Loops can be employed only using the pull semantics, causing data to be processed in nonvectorized manner (in loops).


The second approach promises possibility of unordered solution of pipelines which contain only branching such that all data are processed by SIMD instructions only (except for some epilogues and some split and merge internals). We show that loops can be added to this solution with some memory-consumption penalty without loss of parallel processing. This solution, unfortunately, relies on the ability of deciding whether two conditions are equivalent, which is something we are at the very least not able to verify in some cases since this guarantee may lay in effects. Still, implementation of this solution may be possible with some heuristics on general graphs and is possible with some additional restrictions on form of control flow present in input.


  In the analytical part of this work we implement the mentioned framework. Our implementation is general with respect to SIMD extension, programmememing language and target environment --- it may also be seen as a highly sophisticated text processing environment.


Besides the actual framework, we provide an example instruction table for the Intel Streaming SIMD Extension (SSE) in version 4.2. This instruction set contains all standard operators of the C language implemented for variety of data types (8-64 bit signed and unsigned integers, booleans, floats and doubles). This table, although being thoroughly tested, is not guaranteed to be bug-free due to its vast scope. Also, we implement all graph transformations needed for the ordered solution of control flow. Furthermore, we implement a generic set of macros which implement a register-allocated buffer and use these for implementation of generic split and merge instructions. Using these we implement showcase of the first solution which is capable of handling branching. Unfortunately, our solution has very high overhead on control flow instructions as well as reserves for improvement. 

TODO: be more optimistic

\section{Note on paper structure}

The rest of this thesis is organized as follows:
\begin{itemize}
  \item Chapter \ref{ch:analysis} begins by introduction of formal task assignment. A section which deals with the simple case of pipelines without any control flow follows. Last sections of this chapter discusses the problem of control flow in-depth.
  \item Chapter \ref{ch:implementation} introduces architecture of our framework on high level of abstraction, explains its basic working principles and mentions concrete algorithms which were employed.
  \item Chapter \ref{ch:sse} provides a brief introduction to the SSE extension and explains how the provided example SSE instruction set is built.
  \item Chapter \ref{ch:faq} deals with concrete cases of usage of our framework. Specially we specify here exact semantics of all input interfaces.
  \item TODO: fill in some conclusion.
\end{itemize}




\chapter{Theoretical analysis and solution}

    \label{ch:analysis}
    %\section{Motivation, the bobox environment}

    %define a box
    \section{Formal introduction, definitions, notation} %Reader should be familiar with the basic terminology of the graph theory. We will assume knowledge of the following terms: \emph{directed graph}, \emph{multigraph}, \emph{topological ordering}, \emph{path}, \emph{cycle}, \emph{connected component}, \emph{depth first search (DFS)}. Definitions of these may be found in \cite{kapitoly}. 


We provide a significant number of definitions of our own. We construct these formalisms since this problem has not yet been described in this exact form (as far as we know).  We provide these formalisms mainly for the sake of clarity. We only define the minimum formalism neccessary for this thesis. Since introducing all used formalisms at once does not seem wise, we choose to build our formalisms incrementally. We will provide some forward references and some limited amount of forward context to keep the flow of information as consistent as possible, but our main goal will be providing information at places where they are actually used.


In this section, the reader should get some intuitive understanding of flow graph meaning, of what flow graphs represent and should note the difference between instructions and operations. The reader should feel free to skip formal details.

\subsubsection{Introduction of flow graphs}

\mybegindef{edgeincidence}{Edge incidence}
  Let $G(V,E)$ be a multigraph. We define the relation $\sim$ by the following rules: \begin{itemize}
  \item $\sim \ \subseteq E \times (V \times V)$ 
  \item $ (\forall e \in E)(\forall a,b \in V)( e \sim (a,b) $ if and only if edge $e$ leads from vertex $a$ to vertex $b$ $)$
  \end{itemize}
\myenddef


\mybegindef{powerset}{Powerset}
  We will use the symbol $\powerset(V)$ to denote a \emph{power set} of a set $V$. The power set of the set $V$ is defined as a set of all subsets of the set $V$.
\myenddef


\mybegindef{queue}{Data queue}
  We will use the term \emph{data queue} or simply a \emph{queue} for a first-in first-out container with the following operations:
  \begin{description}
  \item[push] adds an element to the container.
  \item[pop] returns the oldest element present in the container and removes it from the container.
  \end{description}
  We will use this term in its conceptual meaning. In real implementation, we use constant-size buffers or simple memory fields for objects which we theoretically describe as queues.
\myenddef


\mybegindef{flowgraph}{Flow graph}
  We will use the term \emph{flow graph} for any tuple $(G(V,E),O,from:E \arrow \N, to:E \arrow \N, op:V \arrow O,\\ inputs: \arrow \powerset(V), outputs: \arrow \powerset(V) )$ which satisfies the following conditions:
\begin{itemize}
  \item $G(V,E)$ is a nonempty connected directed multigraph with annotations defined by $from$, $to$ and $op$.  
  \item Let $G$ be acyclic except for edges which are explicitly allowed to be on cycles by semantics of $O$ and by node classification provided later. \footnote{This requirement may be ignored for the time being since there are no general restrictions on semantics of $O$. We include this requirement since implicit acyclicality will greatly simplify enumerations of assumptions.}
  \item $O$ is a set of operations.
  \item Vertices in $inputs$ have no incoming edges and vertices in $outputs$ have no outgoing edges. Also, we require that the $op$ annotations of these vertices have semantics of input or output nodes as described later in definition \ref{def:nodetypes0}.
  %\item $\forall(v) \in V |out(v)| \leq 1$ 
  \item Any directed edge is uniquely identified by its destination vertex and the $to$ annotation. I.e.:
    \begin{align*} &(\forall{e,f \in E(G)})(\exists a,b,c,d) \\ &(e \sim (a,b) \land f \sim (c,d) \land b = d \land to(e) = to(f) \Rightarrow e = f )\end{align*}
\end {itemize}
    We will write $G(V,E)$ (with $O$) for brevity.
\myenddef

    This definition resembles the concept of Kahn networks (\cite{kahn}) as well as Hybrid flow graphs (\cite{hfg}).

  Flow graphs represent pipeline computations. Edges represent data flow. The edges may also be understood as data queues. Vertices represent data transformations. The $to$ annotation identifies arguments for operations represented by members of $O$. The $to$ annotation identifies outputs of operations which return more than one value. The nullary functions $input$ and $output$ denote sets of nodes which are meant as data inputs or outputs.

    We provide an example of a simple flow graph with all annotations explicitly shown in figure \ref{fig:flowgraph}. Later, we will omit the $from$ and $to$ annotations entirely while showing the $op$ annotation without explicit labels. 
  \graph{flowgraph}{A flow graph with all annotations shown}

Some semantical meaning of this graph may be:

  \begin{description}
  \item[a,b] - First, the two load (LD) operations produce two values from somewhere, typically from some two (distinct) containers stored in RAM, and store these values into the corresponding outgoing queues of $a$ and $b$ (which will be typically implemented as some register-allocated variables). 
  \item[c] - The DIV operation gathers the two values from the incoming queues, performs division, and pushes the result into its outgoing edge. Note that the $to$ annotation provides the information that the input from $a$ should be divided by the output from $b$ and not vice versa.
  \item[d] - The store operation retrieves the result from its incoming queue and stores them somewhere. Again this operation may represent almost any effect.
  \end{description}

\begin{rem}
  The reader may be asking why is G required to be (implicitly) acyclic. The main reason is that this way of thinking of flow graphs matches the way we use flow graphs. More specifically, the apparent problem is that unrestricted loops may introduce unsolvable dependencies. Later, we define special node types which are allowed to void acyclicality in exchange for a guarantee of not introducing cyclical dependencies. This way, there still exists a topological ordering in its \emph{semantical} sense (although not in the definitorical one). 
  %We alter the definition of a flow graph in section \ref{sec:nodetypes}. 
\end{rem}

\subsubsection{Relation between flow graphs and computations}

  This subsection formalises how flow graphs generally describe algorithms, which roughly corresponds to the information informally introduced by semantics of our example.


    \mybegindef{pipeline}{Pipeline of a flow graph}
  Let $G(V,E)$ with $O$ be a flow graph. We will use the term \emph{pipeline} (of flow graph $G$) for a tuple $P(G(V,E), O, Q, queues)$ which adds a \emph{data queue} to every edge of $G$. I.e., we add a set $Q$ of queues and a bijection $queue: E \arrow Q$. 
\myenddef

This definition adds operational semantics to flow graphs. While flow graphs describe input computations, pipelines will let us reason about actual evaluation.


    \mybegindef{datarow}{Data row}
  Since input data for realisations of flow graphs may be understood as some tables consisting of mutually independent rows of data, we will use the term \emph{data row} for a single set of input values for a flow graph $G$. I.e. we will use the term \emph{data row} for a function $inputs_G \arrow D$ for an arbitrary set $D$ of values which is consistent with semantics of operations of $G$.

  We will also (informally) say that all values computed from a data row belong to this data row. This will be always clear from context.
\myenddef


    \mybegindef{operation}{Operation and instruction}
We will use the term \emph{operation} for general identification of an operation which is to be performed on some data. The term \emph{instruction} will be used for a specific way of performing an operation. This way, multiple instructions may be associated with a single operation.
\myenddef

E.g., multiplication by a constant power of two is a unary operation. Associated instructions may be:
\begin{itemize}
  \item binary shift to left
  \item multiplication using the asterisk operator
  \item a vectorized version of multiplication
\end{itemize}

\mybegindef{operationsemantics}{Semantics of an operation}
  When talking about semantics of an operation (or of a set of operations), we mean:
  \begin{itemize}
    \item Structural requirements on any flow graph which contains vertices representing the operation, such as:
    \begin{itemize}
      \item Requirements on existence and counts of incoming and outgoing edges and their $from$ and $to$ annotations.
      \item Data type semantics of $O$.
      \item Any other semantic restrictions of $O$ relating to the graph structure. 
    \end{itemize}
    \item Description of a computation which transforms its input data into its output data.
  \end{itemize}
\myenddef

  We do not consider consistency of runtime semantics (such as semantics of division by zero) of operations in this text. We assume that either all problematic semantics is solved as part of the semantics defined above or that all data are valid with respect to the structure of the flow graph which is supposed to process them.


\mybegindef{transformation}{Data transformation on vertex}
  Let $P$ be a pipeline of a flow graph $G(V,E)$. \emph{Data transformation} on a vertex $v$ will pop some elements from some incoming queues of $v$ and push some elements into outgoing queue of $v$. An element pushed into an outgoing queue $e$ of $v$ will be determined by the elements taken from the incoming queues of $v$, their $to$ annotations and the annotation $from(e)$ according to the semantics of $op(v)$. The transformation is allowed to perform at most one operation on every input or output queue of $v$.
\myenddef


\mybeginalg{greedy realisation algorithm}{GRA}
  Let $P$ be a pipeline of a flow graph $G(V,E)$ with $O$ with possibly nonempty queues. We shall call the following algorithm \emph{greedy realisation algorithm} (GRA).
  \begin{code}
fun GRA(P as defined above)
{
  while P contains a vertex v s.t. 
    - incoming queues of v contain amount of data 
               required by semantics of op(v) 
    - v does not belong to inputs or outputs of G
  {
    perform transformation denoted by op(v) on P;
  }
  return P;  
}
  \end{code}
\myendalg


\mybegindef{realisation}{Realisation of a flow graph on a single data row}
  Let $G(V,E)$ be a flow graph. \emph{realisation of a flow graph} $G$ will be an algorithm which provided with a single data row $d$ produces content of incoming queues of vertices in $outputs_G$ such that: 
  \begin{itemize}
    \item $G$ is consistent with semantics assigned to O.
    \item Input consists of exactly one value per every outgoing queue of a vertex in $inputs_G$.
    \item Output consists of exactly one value per every incoming queue of a vertex in $outputs_G$.
    \item Output is equivalent to the output obtained by the following procedure and is unique\footnote{This is an implicit requirement on semantics of O. This requirement causes that realisation does not exist if the semantics are not unambiguous.}:
    \begin{enumerate}
      \item Let P be a pipeline of G with queues initialized so that:
      \begin{itemize}
        \item A queue $q \in inputs_G$ is initialized to a single value, defined by the data row $d$.
        \item An empty queue otherwise.
      \end{itemize}
      \item Run GRA on P.
      \item Retrieve all values from incoming queues of vertices in $outputs_G$ as output.
    \end{enumerate}
    \item The GRA is required to consume data from all queues except the output queues, i.e., the incoming queues of $outputs_G$.
  \end{itemize}
\myenddef


\mybegindef{realisationx}{Realisation of a flow graph}
  \emph{realisation of a flow graph} will be an algorithm which provided a flow graph $G$ and multiple data rows produces results equivalent to results obtained by separate application of the \emph{realisation of a flow graph on a single data row} on every data row.
\myenddef


Note that our definition does not perform any input or output operations. The reason for this is that we do not wish to deal with effects in the theoretical analysis.


Also, note that our definition of flow graphs does neither require nor ensure semantics of $O$ to be consistent in G. Also, the definition of a realisation requires but does not ensure that GRA produces exactly one value per every output vertex. Thus, some flow graphs may have no realisation. We will implicitly assume our flow graphs to be consistent with some semantics assigned to $O$. This semantics will always be clear from their context.


\mybegindef{consistency}{Consistency of a flow graph}
  We shall say that a flow graph $G(V,E)$ with a set of operations $O$ is \emph{consistent} if and only if there exists some realisation of this flow graph (on a single data row). 
\myenddef


To summarize the definitions provided so far, we have defined what a flow graph is and how it describes some class of algorithms. We have also described many formal details which may be seen as implementation details but which were crucial for formal correctness of the formalisms provided.

\mybegindef{nodetypes0}{Node types}
  We shall distinguish the following types of operations (this type of operation is to be considered a part of semantics of operation):
\begin{description}
  \item [regular node] will be a node with an operation whose semantics is taking one element from every incoming queue of a vertex and pushing one element into every outgoing queue of this vertex. Data to be pushed into an outgoing queue may depend only on the input data and the \emph{from} annotation of the outgoing queue.
\item [input or output operation] will be nodes which represent either consumption or production of exactly one value per data row.  These will be exactly the nodes contained in $inputs_G$ or $outputs_G$. These will not participate in the GRA algorithm due to reasons described above.
\item [contro flow node] will be any other node. We will classify these in definition \ref{def:nodetypes1} since exact semantics of these nodes is implementation-dependent. 
\end{description}
\myenddef

  Some SIMD related definitions are postponed into section \ref{sec:vecgensec}.




    \section{Code generation} 


        \subsection{Simple code generation} Consider the following problem.

\mybeginprob{simpleproblem}{Simple code generation}
Let $G(V,E)$ be an acyclic flow graph consisting of regular operations. Let $O$ denote some subset of arithmetic operations of the C language. Also, let $G$ be consistent with the semantics of these operations. We wish to generate a C code which realises the graph $G$.
\myendprob

In other words, we wish to generate code for a simple basic block without any control flow, such as the one shown below:

\graph{adder}{Adder computation}

We wish to generate one instruction per every vertex so that every vertex receives data from its predecesors. The resulting fragment of code may look like this:

\mybeginfig
\begin{code}
int var_a = input_0[i];
int var_b = input_1[i];
int var_c = var_a + var_b;
output_0[i] = var_c;
\end{code}
\myendfig{simplecode}{Code fragment realising the adder computation from figure \ref{fig:adder}}

\begin{rem}
  We shall explicitly note that the input and output arrays, as well as the $i$ variable, are some environment-dependent effectfull expressions. These make sense for some loop processal, but are by no means part of a generalised algorithm. Other variable names are chosen arbitrarily by the algorithm. The plus operator is an instruction assigned to the addition operation. Types are part of semantics of operations that produced the results.
\end{rem}

Solution to this problem is straightforward. First, we order vertices of $G$ topologically. Then we take vertices in this order and process them one by one. For every vertex $v$ we use its annotation to transform its input data into its output data. Then we save the output data as a new variable under some unique identifier for later use. 

\begin{samepage}
\mybeginalg{simplegenerator}{Simple code generator}\ 
%\begin{defstyle}
%  Let $G$ and $O$ be defined as in the problem \ref{pro:simpleproblem}. 
%\end{defstyle}
\begin{code}
fun enc(operation o)
  { return any instruction assotiated with o}
fun gen_instruction(instruction, result_name, arg1, arg2, ...) 
  { generate actual text representation from arguments; }
fun vertex::result 
  { return the name used for storage of the result } 
fun vertex::in(index) 
  { return u such that there exists e == (u, v) and to(e) == index; } 

fun generator(graph G)
{
  for( v vertex in V in topological order)
  {
    arguments;
    result_name = "var_$v";
    gen_instruction(enc(op(v), result_name, 
      v.in(1).result(), v.in(2).result(), ...);
    v.result() = result_name;
  }
}
\end{code}
\myendalg
\end{samepage}

\subsubsection{Form of patterns}

We may have noticed that vertices are basically of three types, each of some typical form. 

\begin{itemize}
  \item Input -- \emph{``\textless type\textgreater  \textless variable name\textgreater  = \textless load expression\textgreater ;''}
  \item Output -- \emph{``\textless store expression\textgreater ;''}
  \item Operation -- \emph{``\textless type\textgreater  \textless variable name\textgreater  = \textless operation\textgreater (\textless argument 1\textgreater , \textless argument 2\textgreater );''}
\end{itemize}

This allows us to construct simple patterns which may be used for actual code generation (e.g., by the \emph{gen\_instruction} function shown in algorithm \ref{alg:simplegenerator}). These may be the following:

\begin{itemize}
  \item Input -- \emph{``\$type \$name = \$input;''}
  \item Operation -- \emph{``\$type \$name = \$operation;''}
  \item Output -- \emph{``\$output;''}
\end{itemize}

Where $\$input$ and $\$output$ evaluate to environment-dependent expressions, $\$name$ evaluates to a new identifier and $\$operation$ evaluates to a pattern identified by the $op$ annotation. The operation pattern may contain $\$arg1$, $\$arg2$, $\$arg3$... expressions. These evaluate to the saved names from incoming vertices.

This scheme of generation is actually used --- our generator is driven by a quite complex text-processing system which is based on recursive evaluation of shell-like variables in broader contexts provided by various extensions.


 %describe the dollar notation

        \subsection{Vectorized code generation} \label{sec:vecgensec}

\mybegindef{width}{Instruction width}
We will use the term \emph{width} to denote the number of independent data rows a single instruction operates on when we get to utilisation of SIMD instructions.
\myenddef

\mybegindef{widthconversion}{Width conversion}
  We will use the term \emph{width conversion} for an instruction which takes data from an instruction of some width and transforms them to data suitable for another instruction which is of a different width. 
\myenddef

Consider the following problem.

\mybeginprob{vectorproblem}{Vectorized code generation}
  Let $G(V,E)$ and $O$ be an acyclic flow graph defined as in the \emph{simple code generation} problem. Furthermore, let there be a set $I$ (of SIMD instructions) and a partial function $enc (operation, width)$ \footnote{We are merely saying that we wish to have some vectorized instructions available and that every operation should be defined at least in its nonvectorized version.} $: O \times \N \arrow I$. Let $enc$ be defined at least on $O \times \{1\}$. Also, let there be a set of width conversions sufficient for direct conversion between any two widths present in $I$. Finally, let all widths be powers of 2. We wish to generate C code which will realise the graph G using instructions of highest possible widths. I.e., every vertex will be performed by the broadest instruction its operation is associated with. 
\myendprob

If highest instruction widths of all operations present in g$G$ were the same ones, we could use the \emph{simple generator}. We would simply replace all types by the corresponding vector versions and use vectorized counterparts of all instructions. Typically, the width will vary with the bit length of the data type in question and therefore we cannot use this approach with more than one data type present in $G$.


This algorithm may be easily corrected to the following steps:
\begin{itemize}
  \item Generate a number of SIMD instruction which suffices for processal of some $w$ data rows. Do this for every vertex.
  \item Add some width conversions which adapt data outputted by one instruction for the width of the next instruction.
\end{itemize}

As an example, we provide a simple computation which performs modulo 4 on integers. We use nonvectorized store and load instructions and a vectorized mod4 instruction on 16 bit integers as our instruction set. We represent vectorized instructions by well known C constructs since we believe such representation to be brief and readable.

\graph{mod4}{Computation performing modulo 4 on a data stream.}

\mybeginfig
\begin{code}
uint_16 a_0 = input_a[i++];
uint_16 a_1 = input_a[i++];
uint_32 a_conversion_w2_0 = a_0 | (a_1 << 16);
uint_32 b = a_conversion_w2_0 & 0x00030003;
uint_16 b_conversion_w1_0 = b & 0xFFFF;
uint_16 b_conversion_w1_1 = (b & 0xFFFF0000) >> 16;
output_c[j++] = b_conversion_w1_0;
output_c[j++] = b_conversion_w1_1;
\end{code}
\myendfig{vectorizedfragment}{C fragment realising the modulo 4 computation from figure \ref{fig:mod4} in vectorized manner}

In this example, we have generated two LD instructions (rows 1-2), one MOD4 instruction (row 4) and two ST instructions (rows 7-8). Line 3 merges the outputs of vertex $a$, adapting the data for the vectorized modulo operations. Lines 5 and 6 split its output into two values which are suitable for single-value stores on lines 7 and 8.


This problem may be better visualised by a table whose cells represent actually generated instructions. Numbers in cells represent indices of data rows the instruction processes. The width conversions are not shown, but may come in as new rows. Assume that we want our algorithm to process 4 values in one step instead of 2 (or just assume we are showing two invocations at the same time).

\mybeginfig
\begin{center}
\begin{tabular}{c|c|c|c|c|}
  \cline{2-5}
a (LD) & 1 & 2 & 3 & 4\\
  \cline{2-5}
b (MOD4) & \multicolumn{2}{c|}{1-2} & \multicolumn{2}{c|}{3-4}\\
  \cline{2-5}
c (ST) & 1 & 2 & 3 & 4\\
  \cline{2-5}
\end{tabular}
\end{center}
\myendfig{modfusion}{Table visualising operations joined across different data rows}

\begin{rem}
  The condition requiring all widths to be powers of two ensures two things:
  \begin{itemize}
    \item All data vectors are aligned in the meaning that the difference between indices of the first values of any two instances of instructions is a multiple of GCD of their widths. But GCD of two powers of two is always the smaller of the two numbers. This means that the first index of the wider of two instructions is always aligned with the first index of some instance of the smaller one.\footnote{If the meaning of this claim is unclear, the reader may like to try looking at the table in figure \ref{fig:modfusion} for a few seconds. We have merely said that the left boundary of the 3-4 cell in the MOD4 row is aligned with the left boundary of the cells with index 3 at rows of the ST and LD operations.}
    \item We are able to process the same amount of data ($w$) at every instruction.
  \end{itemize}
  Without this condition, we would need $w$ to be the least common multiple of the maximum widths of all operations in $G$. 
\end{rem}

Thus, pseudo code generating the result from figure \ref{fig:vectorizedfragment} may be:

\begin{code}
for(v in V) in topological order
{
  generate adaptation for width of v
  generate code of v
}
\end{code}

The same algorithm in more detail --- first, we find maximum width g$w$ of an instruction which may be used in $G$ (we do not consider operations which are not in $G$ at all). Then we generate code processing $w$ data rows in one iteration. Again, we take vertices of $G$ in topological order and process them one by one. For a vertex whose broadest instructions is of width $n$ we generate $w/n$ instances of its broadest instruction. Arguments are again retrieved from incoming vertices, only now we work with a set of names (of variables holding  results of the vertex in question) per every width that was generated. I.e., when generating a code of a vertex we look whether there is a set of names (variables) for width $n$. If there is, then we use it. If there is no such set, we use some width conversion to obtain data for width $n$ (i.e. generating a new set of variables for width $n$). We show code of this algorithm in algorithm \ref{alg:vectorgenerator}.

\emph{Exactly} the same algorithm, just with some implementation details filled in:

\mybeginalg{vectorgenerator}{Vectorized code generator}\ \ \ 
\begin{code}
fun gen_instruction(instruction, result_name, arg_1, arg_2, ...) 
  { generate actual text representation from arguments; }
fun gen_conversion_ins(from_width, to_width, result_names, arguments) 
  { generate actual text representation from arguments; }
fun vertex::in(index) 
  { return u such that there exists e == (u, v) and to(e) == index; } 
fun vertex::results(w) 
  { return array of names for width w } 

fun gen_conversion(v, to_width)
{
  from_width = native width of v;
  for(int i = 0; i < w; i+= max(from_width,to_width))
  {
    result_names, arguments;
    for(int j = 0; j < from_width/to_width || j == 0; ++j)
      result_names.push("$v_conversion_w${to_width}_${i+j}");
    for(int j = 0; j < to_width/from_width || j == 0; ++j)
      arguments.push(v.results(from_width)[i+j]);
    gen_conversion_ins(from_width, to_width, result_names, arguments);
    v.results(n).push(result_name);
  }
}
\end{code}
\begin{code}
fun gen_vertex(v)
{
  for(vertex u : incoming to v)
    if( u does not have results for width n )
      gen_conversion(u, n);
  for(int i = 0; i < w/n; ++i)
  {
    result_name = "$v_$i";
    gen_instruction(enc(op(v), n), result_name, 
        v.in(1).results(n)[i], v.in(2).results(n)[i], ...);
    v.results(n).push(result_name);
  }
}

fun generator(graph G)
{
  for(v in G in topological order)
    gen_vertex(v);
}
\end{code}
\myendalg



        \subsection{Indirect width conversions} \label{sec:shortpaths}

One thing which should be adressed at this point is the problem of indirect width conversions. We have avoided this broblem by assuming that direct width conversions always exist in the \emph{vectorized generator} problem (problem \ref{pro:vectorproblem}).  Although this problem does not seem difficult, we have not been able to find an optimal solution which would work in feasible complexity. 

The simples approach is generating all conversions which lay on the shortest (or cheapest) path between the data provider and the data consumer. This approach may produce suboptimal results if there are multiple consumers. This algorithm may be further improved if we recalculate distances after production of every single path (considering vertices of the generated path to have distance 0). Unfortunately, this improvement still does not ensure optimal results, as is shown by the following counterexample.

\graph{pathcounter}{Example showing that iterated shortest path algorithm is not optimal for trees.}

\begin{description}
  \item[Producer, consumers] are standard instructions (representants of operations).
  \item[Other nodes] represent width conversions and are annotated by some cost. These may be replaced by simple paths if we wish to have a counterexample without costs.
  \item[Edges] represent possible data flow, i.e., edges connect nodes whose outputs are of the same width as the input of the other node.
  \end{description}

The improved shortest/cheapest path algorithm may first choose the path $ a \arrow 2 + d \arrow b $ and then the path $ a \arrow 4 \arrow c $, yielding path of cost $6+d$. The cheapest path is apparently $5+d$. The variable $d$ shows that strategies like \emph{cheapest path first} or \emph{most expensive path first} do not solve the problem.

This problem may be formulated as an integer linear programme\footnote{Linear programming formulates problems as systems of linear equations. More information may be found in \cite{ilp}}. However, this does not imply existence of polynomial solution since integer programming is NP-hard. 

\mybeginalg{ilp}{ILP of minimal distribution tree}
Let $P$ be the set of a single producer, $C$ be the set of all consumers and $W$ the set of all possible widths. Let $G(V,E)$ be a graph such that:
\begin{itemize}
  \item $V = P \cup C \cup W$
  \item $E = \{ (u,v) | u,v \in V \land \text{there exists a width conversion from } u \text{ to } v \}$
\end{itemize}
Define variables and constants as follows:
\begin{itemize}
\item $m = \max_{v \in V}deg(v)$
\item $(\forall e \in E)$: let there be a variable $f_e \in \N_0$.
\item $(\forall e \in E)$: let there be a variable $i_e \in {0,1}$.
\end{itemize}
Define restrictions as follows (we construct flow in network):
\begin{itemize}
\item $(\forall v \in W)$: $\sum_{e \in in(v)} f_e - \sum_{e\in out(v)} f_e = 0$ (Kirchhoff's law)
\item $(\forall v \in P)$: $\sum_{e \in in(v)} f_e - \sum_{e\in out(v)} f_e = - | C |$ (source)
\item $(\forall v \in C)$: $\sum_{e \in in(v)} f_e - \sum_{e\in out(v)} f_e = 1$ (sinks)
\item $(\forall e \in E)$: $f_e - m*i_e \leq 0$ (indicator for minimisation)
\end{itemize}
minimise $\sum_{e \in E} i_e$

\begin{proof}[Proof of correctness of a feasible solution]
  Suppose for a contradiction that a feasible solution exists but there is no nonzero flow from $p \in P$ to a vertex $c \in C$. Let $A = \{v \in V | $ path from $p \in P$ to $v$ exists over edges with $ f_e \gt 0\}$. We shall count the sum $S = \sum_{v \in A}(\sum_{e \in in(v)} f_e - \sum_{e \in out(v)} f_e)$ in two ways:
  \begin{itemize}
    \item By sum of the restrictions $S = | C \cap A | - | C |$. This is a negative number since $| A | \lt | C |$.
    \item By contribution of edges:
      \begin{itemize}
        \item Edges in $A$ contribute by $0$.
        \item There are no edges with $f_e \gt 0$ going from $A$ due to the definition of $A$. Thus, all outgoing edges contribute by $0$.
        \item All edges going to $A$ contribute by a non negative number.
      \end{itemize}
      This implies that $S \geq 0$. But that is a contradicts the previous result.
    \end{itemize}
  Thus, if a feasible solution exists, there exists a path from $p \in P$ to any $c \in C$ with nonzero indicators $i_e$. Thus, the solution is correct.
\end{proof}
\begin{proof}[Proof of existence of a feasible solution]
  For every $p \in P$ and $c \in C$ construct a flow of size 1 from $p$ to $c$ . Let $f_e$ be the sum of all these flows. Furtermore, if $f_e \gt 0$ let $i_e$ be $1$. Let $i_e = 0$ otherwise. This evaluation apparently fulfils all restrictions and therefore if a solution of the original problem exists, there exists a feasible solution of the presented ILP.
\end{proof}
\myendalg


    \section{Handling control flow} At this point we would like to introduce control flow into our pipelines (That is branching and loops). Branching may be done in two ways:

\label{sec:branching}
\begin{itemize}
  \item By performing both branches on all data rows and then joining the data rows using a \emph{select} operation. This approach is well described in \cite{secondpaper}. This approach may be easily employed in our framework by regular instructions. 
  \item By reordering data vectors and performing every branch only on data that are supposed to be handled by the branch in question. This approach has large overhead on reordering, but may be worthwhile in case of nested branching.
\end{itemize}

We will investigate only the second approach since the first approach may be easily employed by use of regular operations provided input is in form using these operations and provided that there is an implementation of \emph{select} available. The second approach also allows processal of loops. \footnote{The conventional approach for loops would be expanding the loop body multiple times and then optimizing this result, possibly pipelining operations of one loop iteration over multiple resulting iterations and possibly employing SIMD instructions on nondeterministicaly discovered pairing of instructions\cite{compilers}. Also some more advanced approaches such as polyhedral analysis of loops may be employed\cite{polyhedral}. The conventional methods have the benefit of low overhead and lower independency requirements. }

\label{sec:node_types}

First, we shall introduce some new types of nonregular nodes. These will not be required to fullfill the condition that an operation has to pop/push the same amount of data to/from every queue (defined by regularity in define \ref{def:node_types0}). Note that this definition adresses semantics in context of single data row.\footnote{Semantics in vectorized case are defined to be equal to multiple instances of single context.}

\begin{define}[node types]
  We shall distinguish the following types of operations:
\begin{description}
\item [regular operation] As defined in definition \ref{def:node_types0}.
\item [input or output operation] As defined in definition \ref{def:node_types0}.
  \item [split operation] will take one data input and one condition input (Boolean value typically). Depending on the condition input it will decide upon exactly one of its outgoing $from$ annotations and push data into all outgoing queues with this annotation. All other outgoing queues will remain unchanged. The data input and the condition input will be consumed.
  \item [merge operation] will take two data inputs and one condition input. Depending on the condition input it will take data from one of the inputs and push them to all output queues. Data in the chosen input and in the condition input will be consumed. All other incoming queues will remain unchanged.

  \item [loop merge operation] will take two data inputs. If there is any data in the second input \footnote{This is meant in respect to the definition of relization of a single data row. In real implementation this means that we have to flush entire loop \emph{somehow} before checking this condition.}, this node takes the data from the second input and puts it into all outgoing queues. Otherwise it does the same thing with the first input. Moreover, we will allow the second input to be on cycles, i.e., we will not consider the second input edge when saying \emph{acyclic} or \emph{topological ordering}. By doing this we will not void semantical meaning of flow graphs since standard edges represent data dependencies unlike these edges. 

  %Some generation systems may decide against using the condition input and just take data in unspecified order.
  %The observations would not hold with this.
  %\item [loop operation] will merely stand for combination of other node types. The exact allowed combination is shown below. This combination of nodes will be used as if it were a single node. As a whole it takes two data inputs and two condition inputs and produces two data outputs. The second data input and the second condition input are are allowed to be on cycles.

  \item [loop condition] will be an operation which takes two condition inputs and provides two condition outputs. If there is any data in the second input, this node takes an element from the second input. Otherwise, an element from the first input is taken. Results are shown in the following table (this way one column is applied per every consumed element).
    \begin{center}
      \begin{tabular}{c|c|c|c|c}
        index of input, value & 1, false & 1, true & 2, false & 2, true \\ \hline
        output 1  & true            & true           & false           & false          \\ \hline
        output 2  & true            & false          & true            & false          \\ 
      \end{tabular}
    \end{center}
    This table may be interpretted as follows:
    \begin{center}
      \begin{tabular}{c|p{4cm}|p{4cm}}
        & true            & false \\ \hline
        output 1  & process element which enters the loop & process element which already is in the loop \\ \hline
        output 2  & send this element out of the loop & send this element into next iteration \\ 
      \end{tabular}
    \end{center}
      In addition, we will allow this node to have a non-negative integer as an internal state. This integer will be always equal to the difference of number of true and false values sent into the second output. Furthermore, we will restrict option of taking input from the first input by requirement that (the absolute value of) the difference is less or equal to some positive number. This integer will be denoted by a partial function $limit: V \arrow \N$. The internal state represents the number of data rows present in body of the loop managed by this condition. Since we define semantics of a single data row, this remark may be effectively left out since this condition always holds in an environment containing just one data row (which is how all semantics are defined). 

  \item [explicit queue/explicit buffer] is a queue which takes a single element from its only input and pushes it into all outputs. It will help us formalizing things later. 
\end{description}
\end{define}

Using the standard \emph{split} and \emph{merge} nodes we will be able to perform standard branching, i.e., to realize the \texttt{ if \string{\string} else \string{\string} } construct. The \emph{loop} node construct will allow us to realize the \texttt{ while(condition) \string{\string} } construct. We show this in a picture.

\begin{rem}
  Until now we have used only rectangular nodes in our graphs. From now on we will also use circular/elliptical nodes which will denote connected parts of graphs.
\end{rem}

\label{sec:cf_constructs}
\graph{cf_constructs}
\graph{loop}

Until now we have performed all data transformations \emph{regularly}. Thanks to that we did not need to check consistency of flow graphs (defined in \ref{def:consistency} by properties of realization of graph on a single data row). Now we do.

Apparently not all graphs satisfy the second condition. Here is one such graph. The ADD operation will never process anything since it never gets both operands due to the merge operation.

\graph{inconsistent}

We will not present any way of consistency testing of general flow graphs. We will content ourselves with an observation that nontrivial consistent flow graphs exist.

\begin{observation}[consistency of schemas]
  A flow graph $G$ in form of one of the examples of node type usage is consistent if the following conditions hold:
  \begin{itemize}
    \item Branch or loop bodies, conditions, peamble and epilogue consist only of regular operations. 
    \item There are no concealed paths in between:
      \begin{itemize}
        \item if branch and else branch
        \item any branch and a preamble/condition/epilogue
        \item loop body/condition and a preamble/epilogue
      \end{itemize}
    \item There are no input and output operations contained in loop condition or in loop and branch bodies.
    \item Semantics of $O$ are preserved.
    \end{itemize}
  \begin{proof} 
    Thus, we need to show that provided that every input node produces one value, the GRA algorithm produces exactly one value at every output node and that no value remains in the graph (except for queues of output nodes).

        Assume that there remains a node with some input empty and some input nonempty. 
        \begin{itemize}
          \item All partitions depicted as circular/elliptical nodes consist of regular operations only, hence are acyclic. It holds that either all inputs of such partition receive exactly one input or all receive none. Furthermore it holds, that there exists a path from an input (either input queue of partition in question or an input now) into any queue of any node in such partition. It follows that every operation in such partition is invoked exactly once per every data row, having always exactly one value in all its input, consuming these and producing exactly one value per every output. Thus, all values (except for ones in queues of output operations) are consumed, so the assumed situation cannot take place inside of a single partition.  
          \item We have banned any direct contact between any two partitions which may cause this except for the pair of preamble and epilogue.
          \item Flow of data from every split leads to a merge which merges according the same condition, which means that the two data flows are merged in such manner that quantity of the data rows is preserved (and their order too if there were multiple data rows present in the graph). Thus the branching schema is correct.
          \item Apparently, epilogue in the loop merge schema receives exactly one value unless there is an infinite loop in semantics of $O$, which is a situation which we do not assume in this text.
          \item The loop condition operation in the last schema apparently manages the split and merge nodes in consistent manner, so the epilogue of this schema again receives exactly one value.
        \end{itemize}

    %Assume that We have banned any direct contact between any two queues which may cause this. Thus, the only thing to consider is the merge and the split node. But flow of data from every split leads to a merge which merges according the same condition, which means that the two data flows are merged in such manner that both order and quantity of the data rows is preserved.
  \end{proof}
\end{observation}

\begin{observation}[schema nesting]
  Let $G(V,E)$ with $O$ be a consistent flow graph generated by finite repetition of steps described in observations \ref{TODO} to \ref{TODO}. Let $C$ a connected subgraph of $G$ composed only of regular nodes, i.e., of vertices annotated by regular operations. Now we create a new graph by nesting a graph $H$ corresponding to either of the provided example schemas into $C$. We do so by connecting vertices of $C$ and $H$ by edges and by changing some regular operations into another regular operations. Then the new graph is consistent if the following conditions are satisfied:
  \begin{itemize}
    \item The nested schema satisfies the conditions from previous observation.
    \item The nested schema is connected to C only by the preamble and epilogue.
    \item The nested schema does not contain any input operations. 
    \item We have not introduced any new cycle except for cycles containing second input of any loop node.
    \item Semantics of $O$ are preserved.
  \end{itemize}
  \begin{proof} 
    This follows directly from the proof of the previous observation and the structure of $G$.
  \end{proof}
\end{observation}

\begin{observation}[schema sequencing]
  Let $G(V,E)$, $O$ and $C$ be defined as in the previous observation. Now we take one of the example schemas and embed it multiple times into $C$ in the means of the previous observation. If we use the same condition for all inserted instances of the chosen schema, then we may join the corresponding bodies of the embeded schemas without voiding consistency of the new graph.
  \begin{proof} 
    By structure of $G$ using an argument analogical to the proof of the \emph{consistency of schemas}.
  \end{proof}
\end{observation}


The provided schemas apparently suffice for realization of the standard \texttt{ if \string{\string} else \string{\string} } and \texttt{ while(condition) \string{\string} } constructs. However we will still try to provide solutions for general (consistent) flow graphs. The following example shows a consistent flow graph which cannot be generated by means of the observations. What is of our interest is the fact that consistent merge condition input streams for this example can be constructed using \emph{(pure)} regular operations. It suffices to compute the merge condition before we apply any split.

\graph{diamond}

\begin{description}
  \item[a,b,c] are booleans known in preamble.
  \item[e and f] are streams of predicate \texttt{(a)} split by predicate \texttt{((a\&\&c)||(!a\&\&b))}.
  \item[g] is a stream of predicate \texttt{((a\&\&c)||(!a\&\&b))}.
\end{description}





        \subsection{Partitioning the graph} %At this point we would like to examine possibilities of realization of flow graphs with all types of nodes. 

In the previous section we have shown how flow graphs composed of regular operations may be processed. This approach does not suffice in graphs containing control flow since different nodes may need to process data at different speeds. 

Our strategy is to cut a flow graph into partitions such that every partition consists only of regular operations. Then we add buffers between these partitions. The data will then be processed by a crawler which  goes through our graph and processes its partitions one at a time. Every single partition will be processed by code generated by generators we have already presented. 

First, we need some definitions.

\begin{define}[edge layer]
  Let there be a directed graph (possibly multigraph) $G(V,E)$. We will assign one more numeric annotation to every edge. We will call this annotation \emph{layer}. Thus we have $layer: E \arrow \N_0$. If $layer(e) = n$ we will say that $e$ is on layer $n$. 
\end{define}

Layer will denote type of edge. It will allow us to talk about subsets of edges in a more intuitive manner and also to introduce a more complex graph structure. We show our visual convention in the following picture.

\graph{lines}

\begin{define}[factor graph]
  Let $G(V,E)$ be a directed multigraph and $G_f(V_f \subseteq \powerset(V),E_f)$ be a directed graph. We will say that directed graph $G_f$ is a \emph{factor graph} of $G$ if there exists a function $f: V \arrow V_f$ such that:
  \begin{itemize}
    \item $(\forall u \in V)(\forall c \in V_f)( u \in c \Longleftrightarrow c = f(u)) $
    \item $f(u) = f(v)$ if and only if there exists an undirected path $P_{u,v}$ in $G$ which uses only edges on layer 0.
    \item $(\exists e \in E)(\exists u,v \in V)(e \sim (u,v) \land layer(e) = 1) \Longleftrightarrow \exists(e_f \in E_f)( e_f = (f(u), f(v)) )$  
    \item The edge $e_f$ in the previous condition is unique.
  \end{itemize}
  We will call vertices of $G_f$ \emph{components} since they are connected components in graph defined by $(V, \{e \in E \mid layer(e) = 0 \})$. We will also use the term \emph{component} for partitions of $G$ which correspond to a vertex in the factor graph. I.e., we will call a partition $P$ of $G$ a component if and only if there $(\exists w \in V_f)$ s.t. $(x \in P) \Longleftrightarrow ( f(x) = w)$
\end{define}

\begin{define}
  Let $C$ be a component of a graph $G(V,E)$. We will use the term \emph{input vertex of a component} for a vertex $v \in C$ s.t. $(\exists w \in V)(\exists e \in E)( e \sim (w,v) \land layer(e) \geq 1 )$. We will also use the term \emph{output vertex of a component} defined analogically.
\end{define}

Note that this definition overrides another meaning of this term. Also note that a factor graph exists for every $G$ and that it is unique (up to isomorphism).

Our partitioning consists of the following steps:
\begin{enumerate}
  \item Control flow node expansion. 
  \item Removal of cycles in factor graph. 
\end{enumerate}

Cycles may be introduced into factor graphs by node expansion, even into factor graphs of acyclic flow graphs. Any path in a factor graph represents some path of buffers where data wait for processal. A cycle in a factor graph means that some computation path containing either control flow (lets say $P$) or explicitly inserted buffers leaves a component (lets say $C$) and later joins it. We need to remove all cycles from the factor graph since we want to process components of the resulting graph in a single processing body. We do so by splitting the component $C$ into two new components $C_1$ and $C_2$ such that $C_1 \rightarrow P \rightarrow C_2$.

\subsubsection{Node expansion}

In this step we replace all control flow nodes by new vertices which are connected by layer 1 and 2 edges. By this we insert splits at places where data need to be reordered (where the corresponding streams need to be either split or merged). We show expansions in the following picture.

\graph{expansions}
\graph{expansions_mergeloop}
\graph{expansions_condition}
\graph{expansions_split}
\graph{expansions_queue}

\subsubsection{Cycle removal}

Let $G(V,E)$ be a flow graph. In order to process data we need a factor graph $H$ of $G$ to be acyclic. Unfortunately $H$ may contain cycles and loops even though $G$ (restricted to layer 0 and 1 edges) is acyclic. 


We propose the following algorithm to split components so that cycles get removed.

\begin{samepage}
\begin{code}
  H = G
while ( cycle in H exists )     
{ //loop i
  colour vertices of G black
  C = any component on cycle in a factor graph of H
  for v in all outputs of C     
  { //loop j
    for w s.t. there is path from v to w 
    (using edges on layers 0 and 1)
    {
      colour w red;
    }

    colour v and all nonred inputs of C green;
    for w nonred s.t. there is path from w to a green vertex
    (using edges on layers 0 and 1)
    {
      colour w green;
    }

    for w nonred s.t. there is an undirected path from w to v
    (using edges on layers 0 not containing red vertices)
    { //loop k
      colour w green;
    }

    bool changed = false;
    for edge (w,x) s.t. w is green and x is red
    {
      put and expand an explicit queue on (w,x) in H;
      changed = true;
    }
    if(changed)
      break;
  }
}
\end{code}
\end{samepage}

Ignoring some details this algorithm proceeds as follows:
\begin{enumerate}
  \item Take a queue on cycle (in factor graph).
  \item Colour its predecessors green.
  \item Colour its successors red.
  \item Colour other vertices somehow.
  \item Cut the problematic component along green-red borderline.
  \item Repeat.
\end{enumerate}

\begin{claim}
  Let $G$ be a flow graph consisting of node types defined so far with control flow nodes expanded. Perform the algorithm \ref{TODO} on $G$. We claim that:
  \begin{enumerate}
    \item The algorithm terminates.
    \item The algorithm yields a graph $H$ whose factor graph contains neither loops nor cycles.
  \end{enumerate}
  \begin{proof}
    Let $K$ be a cycle or a loop in factor graph of $G$ and let $C \in K$ ($C$ from algorithm).
    \begin{itemize}
      \item There exists an oriented path $P$ from some $a$ to some $b$ s.t. $a$ and $b$ are two distinct vertices in the same component. This holds because $G$ (restricted to layers 0 and 1) is acyclic.
      \item If we choose $v \in P$ s.t. $v \neq b$ ($v$ from algorithm), then $b$ will be coloured red. Otherwise $a$ will be chosen as $v$ in another iteration since $a$ and $b$ are in the same component. Thus $b$ will be coloured red in some iteration of loop j.  Moreover, $a$ will be coloured green in the same iteration.
      \item Since $a$ and $b$ are in the same component, every path from $a$ to $b$ will be disconnected. Thus, the algorithm works if it terminates. 
    \end{itemize}
      Putting a queue on an edge does not introduce any new path on level 0 between vertices of $G$ in $H$. It also decreases number of paths on level 0 edges in $H$ between vertices of $G$. There is finitely many such paths in $G$. Thus, the algorithm will terminate.
  \end{proof}
\end{claim}

\begin{rem}
  The algorithm will not terminate in case there is a cycle composed of edges on layers 0 and 1. But we have defined node expansion so that all backward edges are on layer 2.
\end{rem}

This algorithm is not optimal. The loop $k$ may be replaced by a minimum cut algorithm, by which the size of the first cut may be minimized. Unfortunately there is no guarantee that such optimization will not increase the minimum costs of the cuts performed later. If we wanted an optimal solution, all cuts would have to be optimized at the same time. 


As an example we show a simple pipeline which performs some arithmetic operations. Reader will surely excuse empty branch bodies (direct edges from split to merge). We begin by a simple representation of a graph which contains a split and merge control flow nodes and one explicitly inserted queue (the BUFFER node).

\graph{partitioning0}

Node expansion produces a new graph with some level 1 edges. This graph already consists of 4 components, but there are some edges which bypass the control flow. These edges introduce a data dependency in the partition 0 which cannot be met without leaving the code dedicated to processing partition 0. This problem may be seen in the included factor graph, which contains an edge going from component 3 to component 0. There is also a loop on partition 0 caused by the edges which bypass the buffer node. 

\graph{partitioning1}

Finally the cycle removal inserts the \emph{auto buffer} nodes, producing a graph with neither cycles nor loops in its factor graph:

\graph{partitioning2}



        \subsection{Order-preserving evaluation} Consider the following problem:

\begin{problem}[ordered realization of control flow]
Let $G(V,E)$ be a consistent flow~graph with $O$ and $I$ defined as in the \emph{vectorized code generation}. Also, let control flow nodes of $G$ be expanded according to the previous section and let the factor graph of $G$ be acyclic. We wish to generate code which will realize graph $G$ on an arbitrary number of data series, utilizing SIMD instructions whenever possible. We also wish the order of data series to be preserved in output.
\end{problem}

\begin{rem} 
Note that by expanding nodes we have obtained a graph which contains \emph{only} regular operations.
\end{rem}

We already know that components of $G$ can be processed as basic blocks. We propose components of $G$ to be connected by constant-sized register-mapped buffers in place of layer 1 and 2 edges. Reason for our decision to make buffers register-mapped is that we wish to minimize communication with RAM. Reason to make our buffers constant-sized is apparently the fact that register space is limited. We will again choose $w$ as a number of data series processed in any partitions. In this proposal we also wish to preserve the order of data series in all merges.

\parspace

To illustrate the problem, we present the following algorithm.

\begin{code}
enter the first compontent, 
  effectively pushing w data records into pipeline.
DFS( component C in factor of G )
{
  if (C has >= w records in every input)
    while (C has >= w records in every input)
      process C;
  else
    abandon current branch;
}
\end{code}

There are two problems:

\begin{itemize}
  \item If we wanted to process all data available in one go, the size of the buffers would have to increase exponentially with every merge, since every merge may have both its buffers full and thus output $2*w$ records.
  \item Since we wish merges to preserve order of all data, there may arise a situation in which a single element will prevent merge. If we waited with processing of this branch until there were $w$ input data series available, the other branch may overfill its buffers.
\end{itemize}

As a solution we propose another algorithm. This algorithm will address the first problem by pulling new data on demand. The second problem will be solved by introduction of pull semantics which will cause a branch to process less than $w$ records if needed. We generate code similar to the following diagram for every component $i$.

\label{ordered_crawler}
\graphn{crawler}

\begin{description}
 \item[partition 0] For simplicity we will assume that all input sources are placed in the partition 0. This is w.l.o.g since we may use a separate counter for every partition. 
  \item[empty inputs] When we say empty input, we mean an input which is supposed to provide data. I.e., on split we do not jump into a branch which is not supposed to provide input.
  \item[terminating variable] When we exhaust all input, we set this variable to zero to indicate that the first partition should process all its remaining data and terminate. Upon termination of any component this variable is incremented and process repeated with the next partition.
  \item[remaining variable] We use this as a variable indicating count of data series which are yet to be retrieved from all data sources. 
  \item[bold line] denotes implicit continuation, i.e., next operation in case condition does not hold. 
  \item[solid line] denotes next operation in case condition holds.
  \item[dotted line] denotes jumps to different partitions (in case condition holds).
  \item[rectangle] A goto label.
  \item[handle termination zero] \ \ \ 
\begin{samepage}
\begin{code}
if( i == 0 && remaining < w )                            
{                                                          
  terminating = 0;                                     
  if(remaining == 0)                                   
  {                                                    
    terminating = 1;                                   
    goto PARTITION_1;                            
  }                                                    
  goto SINGULAR_BODY_0;
}                                                      
\end{code}
\end{samepage}
\item[handle termination nonzero] \ \ \ 
\begin{samepage}
\begin{code}
if ( i == terminating )                                   
{                                                          
  if(an input of i is exhausted)                  
  {                                                        
    terminating++;                                         
    goto PARTITION_i+1;                                
  }                                                        
  goto SINGULAR_PARTITION_i;
}   
\end{code}
\end{samepage}
  \item[i] An index of the current partition. These indices are topologically ordered.
  \item[j] An index of any other partition depending on context. 
\end{description}

\begin{claim}
We claim that the proposed solution works under the stated conditions if all buffers have positive capacity. That is:
\begin{enumerate}
  \item Algorithm produces correct results (if it terminates).
  \item Algorithm does not enter an infinite loop during the pre-termination phase unless there is an infinite loop contained in the semantics of $G$.
  \item Algorithm terminates correctly.
\end{enumerate}
\begin{proof}[Proof of 1]
  Our algorithm never retrieves data from an empty buffer and never pushes data into a full buffer. Also we required $G$ to be consistent with the semantics of control flow nodes, thus all results are correct. 
\end{proof}
\begin{proof}[Proof of 2] So let there be a flow graph $G(V,E)$ and a sequence of data series such that the described algorithm enters an infinite loop which does not process any data. Let $C$ be this loop (subgraph of factor graph $H$ of $G$)\footnote{This is a directed acyclic graph, not a cycle!}. Let $t$ be topological ordering of $H$. Also let there be a description of content of buffers at the moment when algorithm iterates over $C$.
  \begin{itemize}
    \item The asumption that $C \in H$ is correct since if there is an infinite cycle which uses an edge on layer 2 then the algorithm processes data and hence, behaves correctly.
    \item Apparently the algorithm iterates only over singular code of partitions. This means that all partitions on $C$ have either entirely full output or entirely empty input. Also no partition is contained in $C$ more than once.
    \item Let $M = \{v \in C \mid \neg (\exists u \in C)( u <_t v \wedge (u,v) \in C )\}$ (informally all minimal vertices in $C$ with respect to topological order of components restricted to edges in $C$). Let $n$ be the largest index of a data serie such that serie $n$ was processed by all components. Such $n$ exists since every components in $M$ has a full output queue (otherwise our algorithm would dot continue downwards).

    \item Let $a \in C$ be any partition which the $n$th data serie is the last processed one. Then there is an output buffer of $a$ which contains data from the $n$th serie. Let $e \in C$ be an edge which corresponds to this buffer.  Let $b,c \in C$ such that there exists paths $P_{ab}, P_{cb} \in C \wedge e \in P_{ab} \wedge c \in M$. 

    \item Let $m$ be the highest index such that an element from $m$th serie is still present on path $P_{ab}$. 

    \item Partitions $a,c$ have processed the $m$th serie since $a,c \in M$ and $m \geq n$. Partition $b$ has not processed the $m$th serie since an element from the $m$th serie is still on the path $P_{ab}$. That means that there is an element from the $m$th serie on path $P_{cb}$. (This holds since our algorithm never chooses a branch in which no data from the required data serie are to be found, and since the required data serie is exactly the $m$th one. That is also the reason why we did not use the $n$th serie.) But that is a contradiction because one of the two paths is empty.

  \end{itemize}
  \graph{cycle}

\end{proof}
\begin{proof}[Proof of 3]
It suffices to check that the algorithm will process all data of a terminating component and increment the \emph{terminating} variable.
\end{proof}
\end{claim}

\subsubsection{Analysis}
 
TODO: fill something in

TODO: compute buffer sizes




        \subsection{Out-of-order evaluation} Since the cause of ineffectivness of the previous proposal is the requirement stating that data should be processed in order, we will investigate the same problem without this condition. 

\parspace

If we simply dropped the condition, i.e., by letting merges merge in any order, there would arise a problem that data may get reordered differently in different branches. E.g., like in the following diagram.

\graph{ordered_conflict}

This problem may be addressed by ensuring that data series are always handled as one unit, meaning that every split will be performed on all data except data that are not used any further.

\begin{define}[sequential form of control flow]
We will say that a control flow of a flow graph $G(V,E)$ with control-flow nodes expanded and cycles removed is in \emph{sequential form} if the following conditions hold.
\begin{itemize}
  \item All inputs of $G$ are in the same component.
  \item At least one of the following conditions holds for every component $C$ of $G$: 
    \begin{itemize}
      \item $C$ does not have any input from another component of $G$.
      \item All inputs of $C$ are merge nodes such that corresponding inputs of any two input merge nodes of $C$ lead into the same partition of $G$.
      \item All inputs of $C$ are explicit queues leading from another component $D$ of $G$.
    \end{itemize}
  \item At least one of the following conditions holds for every component $C$ of $G$: 
    \begin{itemize}
      \item $C$ does not have any output going into another component of $G$.
      \item All outputs of $C$ are split nodes such that corresponding outputs of any two output split nodes of $C$ lead into the same partition of $G$ and at the same time all these splits use the same condition.
      \item All outputs of $C$ are explicit queues leading into another component $D$ of $G$.
    \end{itemize}
\end{itemize}
\end{define}

\begin{define}[unordered merge] \ \ \ 
\begin{itemize}
  \item \emph{unordered merge operation} will have the same semantic meaning as the standard merge or loop merge operation except it will not take any condition input. This operation may merge streams in an arbitrary but deterministic manner.
\end{itemize}
\end{define}

It is not clear whether every consistent flow graph can be transformed into a flow graph with control flow in sequential form. What is clear is that there exists a nontrivial class of flow graphs which are in this form.

\begin{observation}[schemas with control flow in sequential form] \ \ \ 
  Let $G(V,E)$ be a flow graph s.t. the following conditions hold.
  \begin{itemize}
    \item $G$ is consistent according to \emph{consistency of schemas} (section \ref{sec:cf_constructs}).
    \item $G$ uses only unordered versions of merge nodes in place of standard merge or loop merge nodes.
    \item All input nodes of $G$ are in preamble.
  \end {itemize} 
  Then $G$ may be transformed into a consistent flow graph with control flow in sequential form. This may be achieved by copying the schema in question onto every path between \emph{preamble} and \emph{epilogue} and joining all newly created branch or loop bodies in the means of the \emph{schema sequencing} observation. We may join these by adding any regular operation which does not alter the semantics of $G$.
  \begin{proof} 
  Sequentiality of the resulting control flow follows directly from definitions.
  \end{proof}
\end{observation}

\begin{observation}
  Both \emph{schema nesting} and \emph{schema sequencing} will produce graphs with control flow in sequential form as long as:
  \begin{itemize}
    \item All input graphs have control flow in sequential form and all paths from \emph{preambles} to \emph{epilogues} are transformed in the means of the previous observation.
    \item The condition requiring all inputs to be in the same component (w.r.t. regularity and connectedness) is not voided.
  \end {itemize} 
  \begin{proof} 
  The key observation is that the previous transformation necessarily splits \emph{preamble} from \emph{epilogue} (with respect to regularity of nodes). Thus, nesting of a new schema into a component will never create two paralel schemas except for schemas which comply with the \emph{schema sequencing}.
  \end{proof}
\end{observation}

%\begin{claim}
%Let $G(V,E)$ with $I$ and $O$ be a consistent flow graph with control flow in sequenced form. 
%\end{claim}

Consider the following problem

\begin{problem}[unordered realization of control flow]
Let $G(V,E)$ be a consistent flow~graph with $O$ and $I$ defined as in the \emph{vectorized code generation}. Let $G$ use only unordered versions of merge nodes and let the control flow of $G$ be in sequential form. Let all inputs of $G$ be in one component. Also, let control flow nodes of $G$ be expanded and let the factor graph of $G$ be acyclic. We wish to generate code which will realize graph $G$ on an arbitrary number of data series, utilizing SIMD instructions whenever possible. 
\end{problem}

\begin{claim}
We claim that the \emph{ordered crawler} (section \ref{ordered_crawler}) algorithm solves this problem as well if we make the following adjustments. In addition we claim that this algorithm will always use vectorized version of body until the termination phase is entered.
 \begin{itemize}
   \item All jumps will lead to vectorized labels except for jumps in termination handlers.
   \item If $a$ and $b$ are processing bodies of merge nodes and $C$ is a component such that both $a$ and $b$ outputs data into $C$ then $a$ and $b$ get generated into the same (algorithm) partition.
   \item All buffers have capacity at least $2*w$.
 \end{itemize}
We will again show the following points:
\begin{enumerate}
  \item Algorithm produces correct results (if it terminates).
  \item Algorithm does not enter an infinite loop during the pre-termination phase unless there is an infinite loop contained in the semantics of $G$.
  \item Algorithm terminates correctly.
\end{enumerate}
\begin{proof}[Proof of 1]
  All data series are handled as consistent packs, so every data serie is handled correctly.
\end{proof}
\begin{proof}[Proof of 2] So let there be a flow graph $G(V,E)$ and a sequence of data series such that the described algorithm enters an infinite loop which does not process any data. Let $C$ be this loop. Let $t$ be topological ordering of a factor graph of $G$. 
  \begin{itemize}
    \item Again let $M = \{v \in C \mid \neg (\exists u \in C)( u <_t v \wedge (u,v) \in C)\}$. But $M$ is empty, since every component has either all inputs sufficiently populated or none.
  \end{itemize}
\end{proof}
\begin{proof}[Proof of 3]
This proof has not changed.
\end{proof}
\end{claim}

\begin{rem} 
  The partitioning algorithm may be easily adjusted for conversion of some flow graphs into their sequential equivalents. The simplest version will handle graphs which contain simple edges which bypass some control flow schemas. This algorithm may also be further refined to handle situations like parallel branching. The point of the adustment is in choosing a 'native' outgoing control flow node which will then be pasted on any bypassing edges or before any other control flow which do not match the 'native' control flow node. One question is how to solve some implementation details in order to prevent entering of infinite loops. Another question is how to optimize the resulting graphs, since a naive implementation may produce highly inefficient results.  
\end{rem}




        \subsection{Out-of-order evaluation refinement} The unordered solution seems to be superior to the ordered one in many aspects. An optimal situation would, of course, be if we could return data rows in an arbitrary order. Unfortunately, the data receiver may need the data ordered.  We, of course, may add row indices into the data rows for identification, but that does not solve the problem entirely, since either we will have to store the results one by one or someone else will have to retrieve them one by one. Either version would increase the amount of data transfers between CPU and RAM significantly. 


Ideally, we would like to reorder the data back using as few RAM stores and loads as possible. For this purpose, we may use the information about performed shuffling of the data. We propose the following approach:
\begin{enumerate}
  \item We process all the data and store the results into the RAM. Also, we store a sequence of boolean values per every split or merge node. This sequence will contain a single value per every data row which was processed by the node in question. This value will represent either left or right branch of the split or merge.
  \item We construct a new pipeline which will be an exact copy of the processing pipeline but will be built backwards and will not contain any data processing.
  \item Then we pipe the data from RAM into the new pipeline in inverse order. Also, we provide every split and merge node with the values stored by a corresponding merge or split as a control.
\end{enumerate}

\begin{rem}
  Note that we cannot pipe the original pipeline directly to the one built backwards. Imagine that the first data row to enter the processing pipeline is the last one to come out. If we piped the first pipeline directly into the sorting one, we would have to store \emph{all} the data in CPU registers until the first data row was processed and could be piped into the sorting pipeline.
\end{rem}

This approach surely has its drawbacks, but for some pipelines may reduce the amount of RAM transfers needed for reordering of the data.


\chapter{Architectural overview}

    \label{ch:implementation}

   \section{Brief introduction of components}

Our framework, abbreviated as \emph{ctb}, is implemented as a header library consisting of multiple classes composed via templates. It may be used directly from a C++ code as well as a standalone executable. We understand the data processing to be a sequence of actions which are described by a sequence of commands. 

First, we would like to present all components of our framework on high level of abstraction. We provide more detailed API descriptions in Chapter \ref{ch:faq}.

\subsection*{Command environment}
  The \emph{ctb} class (\ttt{ctb.h}) acts as a frontend of our generator. It contains:
  \begin{itemize}
    \item Methods representing commands.
    \item Interpreter of commands.
    \item Tables containing various access methods of other templated components. These are specifically:
    \begin{itemize}
      \item Table of loaders.
      \item Table of alias environments which act as separate text processors.
      \item Table of alias environments which represent an output environment.
      \item Table of commands.
      \item Table of graph transformations.
    \end{itemize}
    \item Methods which allow registration of new commands and components. This allows custom classes to be used without actual changes to our code.
    \item A single instance of a generator.
    \item A single instance of an instruction table.
  \end{itemize}

\subsection*{Instruction Table} 
  Instruction table (\ttt{instructions.h}) is designed to store definitions of operations, instructions, data types and width conversions. Instruction table supplements the code generator with all information it needs. We provide more detailed semantics at the end of this chapter. For the time being it is more than sufficient to mention that instructions, types and conversions are resolved into string patterns which are later used by the generator. (We refer to these as to \emph{instruction patterns}.)

\subsection*{Graph}

  We provide an implementation of a graph template (\ttt{graph.h}, \ttt{graph\_factor.h}) with the following features:

  \begin{itemize}
    \item Graph may act as both directed and undirected graph.
    \item Vertices and edges are implemented as separate objects. Both vertices and edges contain an instance of a user-defined type. 
    \item Support for distance calculation.
    \item Crawling methods.
    \item Cycle detection.
    \item The template contains data fields for all general annotations specified by our definition of a flow graph. This mainly includes semantics of edge layers and graph factorisation.
    \item Support for computation of a factor graph. The factor is then accessible as a member object.
    \item Export of both the graph and its factor graph into dot representation. 
  \end{itemize}

\subsection*{Generator}
  Generator (\ttt{generator.h}) may be viewed as a special case of the Graph structure. It is parametrised by the type of an instruction table and by the type of a graph. An instance of the generator class generates some instruction code using a provided instruction table (where annotations of graph elements refer to structures contained in the instruction table). 

\subsection*{Loaders}
  Loaders (\ttt{loader\_*.h} files) are classes which provide methods which import and export representations of graphs and instruction tables. These are required to contain a string identifier and four methods which perform the mentioned operations. Every loader is supposed to be registered within the ctb class. The identifier is then used to specify which loader should be used in commands. 

  Exact definition of formats is provided in sections \ref{sec:formats} and \ref{sec:apis}, altogether with the semantics of the generator-related functionality.

\subsection*{Preprocessing environment and the Writer class}

As we have mentioned before, our generator uses a complicated text processing environment to generate code. The key class is the Writer class parametrised by a tree structure of alias environments. 

\subsubsection{Writer}
Writer (found in \ttt{writer.h}) may be simply viewed as containers of generated code or as a string. Writer was designed to act as a data sink for the Generator class. Besides various io and formatting related functionality, the Writer provides a variadic\footnote{Taking a variable number of arguments.} print method. This method takes a format string, which resembles the \ttt{printf}\footnote{C print method which prints formatted text} format but looks and acts more like the text preprocessor of Bourne shell. The print method takes its first argument and recursively performs a sequence of expansions. These expansions may, for instance, perform:
\begin{itemize}
  \item Simple substitution of arguments into the format string.
  \item Evaluation of arithmetic expressions. 
  \item Evaluation of \emph{aliases}. When a variable with a textual name is encountered, it is passed to the root alias environment which is then responsible for resolving of the alias into a result.
\end{itemize}

Expansion formats are further described in Section \ref{sec:preprocessor}.

\subsubsection{Alias environments}
Aliases are string constants standing for other string constants.

Alias environments (\ttt{aliasenv\_*.h}) are \emph{static} classes which define and resolve aliases for writer objects. Alias environments are invoked via a simple method which takes a string (a name of an alias) and returns another string. An alias environment may simply look up the required alias in a map, but may also perform arbitrary computations or even effects. To allow text-processing contexts combining multiple environments, alias classes are supposed to employ some sort of inheritance, therefore creating a tree-like structure for evaluation. 

Alias environments (altogether with the Writer class) represent a generic mechanism which we use in different manners in often unrelated contexts.

For instance, the following alias environments may be found in our framework:
\begin{itemize}
  \item An empty environment.
  \item A bind environment, which takes two other environments as parameters and \emph{combines} them into a single environment.
  \item An aliasenv maker, which is a general environment whose content may be dynamically managed. This is used for purposes which do not have substantial semantics, e.g. various auxiliary usages. 
  \item The control flow macro environment (\ttt{aliasenv\_cfmacros.h}), which implements buffer macro expansions. This is an example of a more involved usage.
  \item The generator alias environment (\ttt{aliasenv\_generator.h}) is used for resolving of generator-related aliases which can be typically found in instruction patterns.
  \item Environments specific for output environments, further described in Section \ref{sec:env2}.
    (\ttt{aliasenv\_simple.h}, \ttt{aliasenv\_cf.h}, \ttt{aliasenv\_simu.h} and \ttt{aliasenv\_bobox.h})). 
\end{itemize}

\subsection*{Alias environments as drivers of the generation}
  Some alias environments also represent output environments --- for instance, a specific language or a specific way how the basic blocks, generated by a generator, should be composed. These environments, registered in an instance of the ctb class, are the actual top-level drivers of the generation. When code generation is invoked, one of these environments is provided with a generator object and is supposed to use the generator (or more accurately its output) to produce some useful code. These environments are required to have a string identifier assigned.

\subsection*{Preprocessors}
  A preprocessor is an alias environment which was registered within an instance of the ctb class as a text preprocessor. A preprocessor may later be used for processing of arbitrary input, meaning that an input file is taken, printed using a writer parametrised by this environment and outputted. This functionality makes sense since nontrivial text transformations may be implemented by means of alias environments.

\subsection*{Transformers}
  A transformer class is a class which provides some graph transformation. This allows arbitrary graph transformations to take place during various phases of computation of our framework. Transformation classes may be registered and used in the same manner as loaders. Again, transformers are identified by string names.

  Currently, the only transformer is to be found in file \ttt{cf\_transofm.h}.

\subsection*{Other structures}

\subsubsection{Cartesian multiplier}
   Cartesian multiplier (\ttt{cartesian\_multiplier.h}) is a simple (but yet interesting) utility which helps to generate cartesian products of containers. The cartesian multiplier is basically an iterator which dereferences to a list of iterators. Values of iterators in this list change with every iteration of the top-level iterator so that all combinations may be observed.

\subsubsection{Parser}
  The Parser class (\ttt{parser.h}) is a simple push down automaton which evaluates arithmetic expressions. This parser was handwritten but its working principle resembles standard LR parsers. \footnote{More on the topic of languages and automatons may be found in \cite{automatons}.}

\subsubsection{Languages}
  Besides other features, the Writer class provides code formatting and indentation. Language classes (\ttt{languages.h}) contain sets of callbacks which are used to determine indentation and line breaking. These are passed to an instance of the Writer class as member typedefs of alias environments.

\subsubsection{Implicit containers}
  There are two container templates (in \ttt{multicont.h}) which basically add one dimension to the template parameter. We use these at places where we needed to replace a single (always present) element by a list of optional elements without invalidation of already existing code. These also serve as a syntactic sugar, since we do not require any container semantics unless the user explicitly needs to use multiple instances of the object in question.

  Neglecting the fact that \ttt{int} may not be used in place of a class, implicit container of integers could work as follows:
\mybeginfig
\begin{code}
impl_contA<int> a = 3;
int b = a + 5; //8
a[3] = 4;
int c = a + a[3]; // 7
int d = a[0] + a[3]; // 7

impl_contB<int> a = 3;
a["c"] = 1;
int e = a + a["default"] + a["c"]; //7
\end{code}
\myendfig{implcontA}{An example of use of implicit containers.}

\subsubsection{Tag handlers}
  Among other characteristics, any instruction may have a list of tags specified. These tags are used to decide whether an instruction may be used or not. The tag handlers (\ttt{taghandler.h}, \ttt{tagmaster.h}) are objects decide of feasibility of a supplemented list of tags. In other words, these objects provide a single, arbitrary function from a set of strings to Boolean.

  Our implementation of this structure provides more involved features and serves for translation of user-provided flags into masks which are later used for determining whether an instruction may be used.

\subsubsection{Proxies}
  During the initial phase of development, we used to use the Proxy class (\ttt{proxy.h}) to make some class members read-only for public. This seemed to be a conceptually good idea, until the management of templated inheritance turned this concept into a nightmare. 



    
In this section, we introduce the employed graph-related algorithms. 

We provide code samples in a partially functional pseudo language inspired by C++11. We omit all unnecessary keywords, signatures and parts of signatures which are not necessary for semantical understanding. We do so to simplify and shorten the code.

Code samples shown in this chapter aim to explain the main ideas without all implementation details. The actual implementation may be found in files \ttt{graph.h}, \ttt{graph\_factor.h} and \ttt{cf\_transformer.h}.

\algref{Graph crawler}
We implement most of our algorithms using a relatively simple crawling template. This template is designed to be usable for wide variety of tasks. Its semantics roughly corresponds to the following pseudo-code (exact implementation may be found at the very end of the \ttt{graph.h} header file). This template was designed to crawl graphs in topological order.

\mybeginfig %
\begin{code}
fun crawler<recurse, inverse>(node* n, fun f, fun g, list l = {0})
{
  if(!g(n))
    return
  if(recurse)
    foreach( u : parent of n)
      crawler(u, f, g)
  if(f(node))
    foreach( u : child of n)
      crawler(u, f, g)
}
\end{code}
\myendfig{crawlerapi}{Our callback-driven crawling template.} %

\begin{description}
  \item \ttt{l} list of edge layers to be used.
  \item \ttt{Child (parent)} of vertex v is a vertex u such that there exists an edge (v, u) ( (u,v) ).
  \item \ttt{f,g} are possibly effectful functions of type $V \arrow \mf{Boolean}$. The function \ttt{f} is designed to perform some work while the \ttt{g} function is designed to control which nodes are to be crawled.
  \item \ttt{Inverse} switches direction of this algorithm, i.e., makes it consider edges in an inversed direction.
  \item \ttt{Recurse} turns on and off the search of parents.
\end{description}

We implement the second recursive call using an out-of-scope queue which is passed implicitly to prevent some loops on call stack. This precaution reduces the maximum number of nested calls to the longest directed \emph{stroke}\footnote{A path which may contain vertices multiple times.} which is present. This functionality (if used properly) allows algorithms to be written so that the maximum number of nested calls does not exceed the length of the longest directed \emph{path}. The actual implementation also handles crawling on a defined subset of graph layers passed by the argument \ttt{l}.

\algref{Topological search}
We implement topological search of connected partitions using the following snippet. This function is very useful since it stands for \ttt{foreach(v in topological order) f(v);}. If there is no topological ordering for the component in question, then this code is equivalent to an unordered \ttt{foreach}.

\mybeginfig %
\begin{code}
fun crawl_partition_topological(node* n, fun f)
{
  crawl<true, inverted>( n
    , (node* n){ f(n); n->lastpass = passid; return true;}
    , (node* n){ return n->lastpass != passid;}
  );
}
fun crawl_topological(fun f)
{
  passid++;
  foreach(vertex v)
    if(v->lastpass != passid)
      crawl_partition_topological(v, f);
}
\end{code}
\myendfig{crawlertopo}{Topological search algorithm.} %

There are only three things to note:
\begin{itemize}
  \item The \ttt{lastpass} variable is used in combination with the (out-of-scope) \ttt{passid} variable as an auxiliary mark which denotes whether a node has been visited. The \ttt{passid} ensures that a new mark which has not yet been used is used.
  \item This function is not re-entrant as shown. In real implementation, we use a bit more complex call which uses an out-of-scope set instead of the \ttt{passid} variable. (This change was made for user convenience when we made the API public.)
  \item We could write this in an even simpler manner, using only the \ttt{g} callback. Unfortunately, such (simplified) solution would enter an infinite loop if there was a cycle in the presumably acyclic graph. The shown solution does not enter an infinite loop, which is the preferred behaviour. 
\end{itemize}

\algref{Cycle detection}
We detect cycles in two steps:
\begin{enumerate}
  \item We construct some ordering of nodes using the topological search.
  \item We check that the ordering is a correct topological ordering.
\end{enumerate}

Both these steps may be done by simple calls to the topological search algorithm.

\mybeginfig %
\begin{code}
fun find_cycle()
{
  int order = 0;
  node* node_on_cycle = NULL;
  crawl_topological((node* n){n->order = order++;});
  crawl_topological((node* n){
    forall(e output edge of n) 
      if(e->from->order > e->to->order)
        node_on_cycle = n;
  };
  return node_on_cycle;
}
\end{code}
\myendfig{findcycleapi}{Cycle detection algorithm.} %

\algref{Shortest path search}
For the purpose of shortest-path routeing, we maintain a map of size $| G |$ in every vertex. This map contains one record per every vertex of $G$. This record contains:

\begin{itemize}
  \item distance
  \item next node on a shortest path
\end{itemize}

This solution, despite being quite minimalist, requires $O(n^2)$ memory, which may be a problem for big graphs. Since we use minimal path search only for the purpose of finding shortest width-conversion path, this limitation does not concern us much. An alternative would be running Dijkstra's algorithm on every request \cite{chapters}.

We fill these routing structures using the well known Bellman-Ford Algorithm \cite{chapters} implemented using our callback machinery. Equivalent pseudo code may be:
\mybeginfig %
\begin{code}
foreach(vertex v)
{
  crawl<false,false>( v
    , (node* n){ recalculate n; return updated || first pass;}
    , (node* n){ return true }
  ); 
}
\end{code}
\myendfig{bellmanfordapi}{Pseudo code of the Bellman-Ford algorithm.} %

\algref{Factor graph construction}
Construction of a factor graph is again a simple process. The only interesting part is the second step in which we try to sort the components (pseudo) topologically even if no topological ordering exists. We do so for aesthetic reasons.

%\footnote{We admit, that we do not have any reason for doing so except for our aesthetical feeling.}.

We should remind the reader that we use two instances of the standard graph to represent the graph and its factor graph. The actual implementation just uses inheritance to allow the factor graph to be a member of the graph (\ttt{class G: public graph\_t<...> \{ graph\_t<...> factor; \};}). For now, we will simply use the names \ttt{G} and \ttt{factorgraph}.
\begin{enumerate}
  \item The crawling algorithm is used to assign component ids\footnote{We use the term \emph{class id} in the actual implementation. This is in accordance with standard algebraic terminology.} to all vertices. We show this in Figure \ref{fig:factoralgclasses}.
  \item In the next step, we try to remap component ids so that the components are in an order resembling our subjective feeling of a correct (pseudo) topological ordering. 
    \begin{itemize}
      \item This task is simple if topological ordering exists - a single pass by the topological search in the factor graph yields a map which is then used to reassign the ids. We show this in Figure \ref{fig:factoralgorder1}
      \item If no topological ordering of components exists, we still try to return nicely-looking results. This makes sense since the underlying graph will usually have a topological ordering. 
        
        First, we assign some ordering to \emph{edges} on levels 0 and 1. We do this by the topological search, so we assume to get the edges ordered (mostly) topologically\footnote{Note that our use case is exactly that of graphs which are acyclic on levels 0 and 1.}. Our observation is that some feasible ordering of components is an ordering induced by the order of input edges. Specifically, we pick the  maximum input edge which is less than every output edge of a component. We show this in Figure \ref{fig:factoralgorder2}.
  \end{itemize}
  \item Finally, the vertices are gathered into components represented by vertices in the new graph structure (as shown in Figure \ref{fig:factoralgtranslate}).
\end{enumerate}

\FloatBarrier

\mybeginfigloose
\begin{code}
crawl_topological((node* n){n->classid = -1;});
int classcount = 0;
foreach(node n in G)
  if(n->classid == -1)
  {
    crawl_partition_topological(n, (node* n){n->classid=classcount;});
    classcount++;
  }
\end{code}
\myendfig{factoralgclasses}{Factorisation -- class assignment.}

\mybeginfigloose
\begin{code}
map<int, int> translation;
int i = 0;
factorgraph->crawl_topological((node* n){
  translation[n->classid] = i++;
});
foreach(node n in G)
  n->classid = translation[n->classid];
\end{code}
\myendfig{factoralgorder1}{Factorisation -- construction of a translation map for acyclic graphs.}

\mybeginfigloose
\begin{code}
for( i in 0 .. classcount-1)
  create component i in factograph (i.e. a standard vertex)
foreach(node n in G)
  assign n to component n->classid of factorgraph;
\end{code}
\myendfig{factoralgtranslate}{Creation of the factor graph itself.}
 
\mybeginfigloose
\begin{code}
int order = 0;
factorgraph->crawl_topological(
  (node* n){ 
    foreach(e output edge of n) 
      e->order = order++;
  }
);
map<int, int> minouts;
map<int, int> maxins;
for(i in 0 .. classcount-1) 
{
  //this ensures that minouts values are always unique
  minouts[i] = MAXINT - classcount + i;
  maxins[i]= -1;
} 
factorgraph->crawl_topological(
  (node* n){
    foreach(e is output edge of n)
      minouts[n->classid] = min(e->order, minouts[n->classid]);
    foreach(e is input edge of n)
      maxins[n->classid] = max(e->order, maxins[n->classid]);
  }
); 
for(i in 0 .. classcount-1)
  maxins[i] = min(maxins[i], minouts[i]);
vector<int> translation = keys of maxins sorted by value;
foreach(node n in G)
  n->classid = translation[n->classid];
\end{code}
\myendfig{factoralgorder2}{Factorisation -- construction of a translation map for possibly cyclic graphs.}


\FloatBarrier

\algref{Buffer size assignment}

This algorithm is implemented according to Algorithm \ref{alg:balancer}. 

\mybeginfig
\begin{code}
map<classid, int> order;
G.factor.crawl_topological(
  (node* p){
    foreach(edge e : input of p)
      order[p->id] = max(order[e->from->id]+1, order[p->id]);
  }
);
G.crawl_topological(
  (node* n){
    foreach(e : output layer 1 edge of n ) {
      coef = order[e->to->classid] - order[e->from->classid];
      e->from->data.set_param("buffercoefout"+(e->from_pos), coef);
      e->to->data.set_param("buffercoefin"+(e->to_pos), coef);
    }
  }
);
\end{code}
\myendfig{balancer}{Buffer size assignment algorithm.}

The mechanism of setting buffer sizes via parameters is further described in Section \ref{sec:splitsmerges}.

\algref{Cycle removal}

We implement this algorithm as described by Algorithm \ref{alg:cycleremalg}. We believe that the reader can imagine how our crawlers can be used to perform the mathematically described steps.

\begin{rem}
  The graph transformer provides a step-by-step visualisation of this transformation. It may be invoked by an additional \ttt{show} parameter passed to the transformation command.
\end{rem}


\section{Instruction representation}

\label{sec:instab}

We would like to shed some light on the structure of the instruction table to make the internal workings of the generator clearer. These structures may be found in (\ttt{instructions.h}). We shall postpone exact field semantics into Chapter \ref{sec:formats}, which is dedicated to API description. The architecture may be visualised by the following diagram:

\graph{instab}{Architectural diagram of the instruction table.}

\begin{description}
  \item [Nodes] of this graph correspond to either \ttt{struct}s or \ttt{class}es.
  \item [Leaf node] structures represent actual \emph{instruction patterns} except for expansions.
  \item [Edges] may be read as `a manages b' or `a contains b' or `a is responsible for b'.
\end{description}

\begin{description}
  \item[Instruction table] contains:
    \begin{itemize}
      \item Containers of \emph{types}, \emph{operations}.
      \item An accumulated search table of \emph{expansions}.
      \item An API which provides mapping from identifiers to these structures.
      \item \emph{Tag handlers} which are used to restrict the set of employed instructions and their access API.
    \end{itemize}
  \item[Operation] class represents either logical operations or expansions. In our implementation every operation is allowed to have at most one output. This is w.l.o.g.. Operations contain/provide:
    \begin{itemize}
      \item An \emph{identifier}.
      \item A container of \emph{instructions} realising this operation.
      \item A container of \emph{expansions} which are identified by operation id of the operation.
      \item Operation semantics such as \emph{width}, a list of \emph{input types} and a list of \emph{output types}.
      \item Publicly available references to \emph{return types}.
      \item An API for retrieval of \emph{instruction patterns}. The generator communicates directly with an instance of the operation class, asking for results suitable for a specific width. The operation then returns patterns of the most suitable instruction which realises the operation.
    \end{itemize} 
  \item[Instruction] is a small structure which represents a target instruction. This structure contains:
    \begin{itemize}
      \item \emph{Width} information.
      \item A list of \emph{tags}.
      \item An \emph{indicator} which indicates whether the instruction is suitable for the current pass of the generator.
      \item \emph{Instruction patterns}. This encompasses some predefined, specially handled fields and a map of user-named fields, which get outputted into different layers of the supplemented writer\footnote{This is implemented by means of one of our implicit containers.}.
    \end{itemize}
  \item[Expansion] may be seen as a special case of an operation or as a placeholder which provides some information about what should be done with some part of a graph. All expansion nodes are supposed to be removed by some graph transformation prior to the generation and possibly replaced by new nodes or even by new subgraphs. The motivation for these nodes are exactly graph transformations, due to clash of the following two requirements:
    \begin{itemize}
      \item When generating code, we need to work only with specific type-aware instructions. 
      \item The graph transformations and node types which we have introduced are type independent. 
    \end{itemize}
    These expansions present a way of binding specific instructions to general transformations by providing a list of instruction identifiers identifying specific target instructions for specific type signatures. These expansions are hence allowed (and assumed) to be present in multiple instances with different type signatures. These contain the following fields:
    \begin{itemize}
      \item A \emph{transformation name} and a \emph{transformer identifier}, which define the graph transformer that should process them and also how it  should process them. We should note that graph transformers are free to ignore these identifiers.
      \item A \emph{transformation identifier}, which also acts as the operation id.
      \item A \emph{type signature}, which is used to choose specific expansion from the set of expansions identified by the same identifier.\footnote{This mechanism allows very simple introduction of type inference. This would be implemented as a graph transformation and its full function would still remain under control of the user via definition of expansions.}
      \item A list of \emph{instruction identifiers} which specify some target instructions for the expansion.
    \end{itemize}
    For instance, consider explicit buffers, defined in Definition \ref{def:nodetypes1}. Buffers may be represented by single nodes in the input graph. This identifier may have two expansions defined. When the `cf' transformer is run, it should choose one of the expansions depending on types of neighbouring nodes and expand it to two new nodes.
    \mybeginfigindent
    \begin{center}
      \begin{tabular}{c|c|c|c|c|p{4cm}}
        identifier & inputs & outputs & name & transformer & list of inst.\\
        \hline
        BUFF & INT & INT & buffer & cf & BUFF\_ST\_INT, BUFF\_LD\_INT \\
        BUFF & BOOL & BOOL & buffer & cf & BUFF\_ST\_BOOL, BUFF\_LD\_BOOL 
      \end{tabular}
    \end{center}
    \myendfig{bufferexpansion}{An example of expansion semantics.}
  \item[Type] represents logical data type of the target architecture. It contains:
    \begin{itemize}
      \item Container of \emph{type versions}.
      \item Container of \emph{type conversions}.
      \item \emph{API} which resolves requests for instruction patterns of type versions and conversions.
    \end{itemize}
  \item[Type version] is a variation of a data type specific for a width. It contains:
    \begin{itemize}
      \item \emph{Instruction pattern} representation of this type in the target language.
      \item \emph{Width}.
    \end{itemize}
    For instance, the type \emph{bool} may have \emph{bool/1}, \emph{int/16}, \emph{int/32} and \emph{\_m128i/128}  as its type versions. (\emph{type pattern/width})
  \item[Width conversion] is a structure which represents a width conversion between two widths of a data type. These contain:
    \begin{itemize}
      \item \emph{Input width} and \emph{Output width}.
      \item \emph{Tags}.
      \item \emph{Flag} indicating whether tag requirements are satisfied.
      \item \emph{Instruction patterns}.
    \end{itemize}
\end{description}

Consider again code fragment given in Figure \ref{fig:vectorisedfragment} (fragment of the vectorised generator). Lines 1, 2, 4, 7, 8 were generated from instruction patterns of \emph{instructions}. The lines 3, 5 and 6 were generated from instruction patterns of two \emph{width conversions} (the lines 5 and 6 are two patterns defined by a common width conversion). All type specifiers present in the figure were resolved from input and output widths and types of instructions and of width conversions. The patterns of instructions and the types were composed together by internal mechanics of the writer. (However, the form of this composition was defined by the form of the pattern. The writer itself is an independent text processor which has \emph{no} information about internal needs of the generator.)

The example in Figure \ref{fig:expeval} shows how a single line of the resulting code may be composed from various patterns. We show a single call to the writer as it may be issued by the generator. Some context is provided via string declarations. The real situation is more complicated since this example does not consider vectorisation. Assume the target environment to be represented by a fictive `clang' alias environment which inherits the `generator' alias environment. 

This example illustrates the context of the generator, of instruction patterns and of instruction-related structures. However, the main point of this example lays in the principle of the text processing environment, e.g., of cooperation between a writer and alias environments.  Understanding of the principle of alias environments is crucial for understanding the entire system.


\mybeginfig
\begin{loosecode}
//determined by the node type
string pattern_type = "$opexpr";
//retrieved from the instruction table
string type_code = "int";
string ins_code = "$arg1 + $arg2"; 
//memoized names of variables containing 
//  results of the operations of vertices a and b
string arg1 = "var_a";
string arg2 = "var_b";
//a new variable name constructed from id of vertex c
string name = "var_c";
//the actual call
w.print("$opexpr", type_code, ins_code, width, arg1, arg2, name);
//the process of evaluation
//    ("$opexpr" -> "$type $name = $operation;") //by clang aliasenv
//"$type $name = $operation;"                //intermediate result
//    ("$type" -> "$1")                      //by generator aliasenv
//    ("$1" -> type_code)                    //by writer
//"int $name = $operation;"
//    ("$name" -> "$6")                      //by generator aliasenv
//    ("$6" -> name)                         //by writer
//"int var_c = $arg1 + $arg2;"
//    ("$arg1" -> "$4")                      //by generator aliasenv
//    ("$4" -> arg1)                         //by writer
//"int var_c = var_a + $arg2;"
//    ("$arg2" -> "$5")                      //by generator aliasenv
//    ("$5" -> arg2)                         //by writer
//"int var_c = var_a + var_b;"               //the final result
\end{loosecode}
\myendfig{expeval}{Example of instruction pattern evaluation.}


\begin{description}
  \item \ttt{type\_code, ins\_code, pattern\_type, arg1, arg2} -- some string variables which were retrieved from contextual information, i.e., either directly or indirectly from the structure of the processed graph or from instruction tables
  \item \ttt{w.print...} -- the actual code which prints an instruction pattern into the writer which acts as data sink for the generator
  \item \ttt{comments} -- below the call, the comments show the intermediate results of the writer and the flow of evaluation
\end{description}






    \section{Actual implementation of the generator}

The generator is implemented according to the scheme described by algorithm \ref{alg:vectorgenerator}. Besides that, the generator provides the writer with a consistent context retrieved from the input graph and the instruction table as described in section \ref{sec:instab}. All other functionality lays either directly or indirectly in alias environments and instruction patterns. (Or is an implementation detail)

We should explain two more things. The first is related to the api of the generator. The generator is directly parametrized by a multilayer writer (created by means of one of the previously described implicit containers) and also by a configuration which specifies how the layers should behave. There are options which allow some layers to receive input only from the first instance of an instruction or receive only the first string written. This allows the generator to output an arbitrary number of code fragments per every single basic block. The former option may be used for simple accumulation of conditions related to control flow semantics of control flow nodes. The latter may be used for indication of presence of specific instructions of for indication of conditions whose value is not operation-specific.

The second thing which is yet to be discussed is the problem of indirect width conversions. We have avoided this broblem by assuming that direct width conversions always exist in the \emph{vectorized generator} problem (\ref{pro:vectorproblem}).  Although this problem does not seem difficult, we have not been able to find an optimal solution which would work in feasible complexity. 

Our implementation deals with the problem by generating all conversions which lay on the shortest (or cheapest) path between the data provider and the data consumer. This approach may produce suboptimal results if there are multiple consumers. This algorithm may be further improved if we recalculate distances after production of every single path (considering vertices of the generated path to have distance 0). Unfortunately, this improvement still does not ensure optimal results, as is shown by the following counterexample.

\graph{pathcounter}{Example showing that iterated shortest path algorithm is not optimal for trees.}

\begin{description}
  \item[Producer, consumers] are standard instructions (representants of operations).
  \item[Other nodes] represent width conversions and are annotated by some cost. These may be replaced by simple paths if we wish to have a counterexample without costs.
  \item[Edges] represent possible data flow, i.e., edges connect nodes whose outputs are of the same width as the input of the other node.
  \end{description}

The improved shortest/cheapest path algorithm may first choose the path $ a \arrow 2 + d \arrow b $ and then the path $ a \arrow 4 \arrow c $, yielding path of cost $6+d$. The cheapest path is apparently $5+d$. The variable $d$ shows that strategies like \emph{cheapest path first} or \emph{most expensive path first} do not solve the problem.

This problem may be formulated as an integer linear programme\footnote{Linear programming formulates problems as systems of linear equations. More information may be found in \cite{TODO}}. However, this does not imply existence of polynomial solution since integer programming is NP-hard. 

\mybeginalg{ilp}{ILP of minimal distribution tree}
Let $P$ be the set of a single producer, $C$ be the set of all consumers and $W$ the set of all possible widths. Let $G(V,E)$ be a graph such that:
\begin{itemize}
  \item $V = P \cup C \cup W$
  \item $E = \{ (u,v) | u,v \in V \land \text{there exists a width conversion from } u \text{ to } v \}$
\end{itemize}
Define variables and constants as follows:
\begin{itemize}
\item $m = \max_{v \in V}deg(v)$
\item $(\forall e \in E)$: let there be a variable $f_e \in \N_0$.
\item $(\forall e \in E)$: let there be a variable $i_e \in {0,1}$.
\end{itemize}
Define restrictions as follows (we construct flow in network):
\begin{itemize}
\item $(\forall v \in W)$: $\sum_{e \in in(v)} f_e - \sum_{e\in out(v)} f_e = 0$ (Kirchhoff's law)
\item $(\forall v \in P)$: $\sum_{e \in in(v)} f_e - \sum_{e\in out(v)} f_e = - | C |$ (source)
\item $(\forall v \in C)$: $\sum_{e \in in(v)} f_e - \sum_{e\in out(v)} f_e = 1$ (sinks)
\item $(\forall e \in E)$: $f_e - m*i_e \leq 0$ (indicator for minimisation)
\end{itemize}
minimise $\sum_{e \in E} i_e$

\begin{proof}[Proof of correctness of a feasible solution]
  Suppose for a contradiction that a feasible solution exists but there is no nonzero flow from $p \in P$ to a vertex $c \in C$. Let $A = \{v \in V | $ path from $p \in P$ to $v$ exists over edges with $ f_e \gt 0\}$. We shall count the sum $S = \sum_{v \in A}(\sum_{e \in in(v)} f_e - \sum_{e \in out(v)} f_e)$ in two ways:
  \begin{itemize}
    \item By sum of the restrictions $S = | C \cap A | - | C |$. This is a negative number since $| A | \lt | C |$.
    \item By contribution of edges:
      \begin{itemize}
        \item Edges in $A$ contribute by $0$.
        \item There are no edges with $f_e \gt 0$ going from $A$ due to the definition of $A$. Thus, all outgoing edges contribute by $0$.
        \item All edges going to $A$ contribute by a non negative number.
      \end{itemize}
      This implies that $S \geq 0$. But that is a contradicts the previous result.
    \end{itemize}
  Thus, if a feasible solution exists, there exists a path from $p \in P$ to any $c \in C$ with nonzero indicators $i_e$. Thus, the solution is correct.
\end{proof}
\begin{proof}[Proof of existence of a feasible solution]
  For every $p \in P$ and $c \in C$ construct a flow of size 1 from $p$ to $c$ . Let $f_e$ be the sum of all these flows. Furtermore, if $f_e \gt 0$ let $i_e$ be $1$. Let $i_e = 0$ otherwise. This evaluation apparently fulfils all restrictions and therefore if a solution of the original problem exists, there exists a feasible solution of the presented ILP.
\end{proof}
\myendalg



\section{Summary}

We have already described all significant functionality of our framework. The next chapters describe only implementation details. The reader should now think over all introduced concepts and the architecture of our framework. Indeed, we have designed a framework which can perform all needed steps to output results equivalent to the algorithm \ref{alg:crawler}:
\begin{itemize}
  \item Our framework can generate vectorized basic blocks as proposed by algorithm \ref{alg:vectorgenerator}.
  \item The framework can parititon the graph as described in section \ref{sec:partitioning}:
  \begin{itemize}
    \item Node expansion may be performed by use of a graph transformer whose input is described by expansions.
    \item Cycle removal is implemented according to algorithm \ref{alg:cycleremalg} by means of a graph transformer.
  \end{itemize}
  \item Code of a partition, as described by algorithm \ref{alg:crawler}, may be generated by a custom alias environment. This environment would contain a static template of this partition with placeholders in place of all conditional jumps specified by the algorithm. As proposed, the control flow nodes may accumulate their control-flow related conditions into named layers of the writer. Substituting these accumulated values into the template is then a simple procedure consisting of a single call to a new writer.
  \item Other generic procedures, such as buffers may be implemented either externally in the C preprocessor or by means of alias environments. Either implementation is the interfaced via records in an instruction table. We employ the latter way due to stubborn recursion-rejectiveness of the C preprocessor\footnote{Recursion may actually be achieved by means of various workarounds. Unforunately, the implemented algorithms tend to fail due to obscure reasons.}. Specifically, we implement an in-register SSE buffer in the the `cf\_macros' environment.
\end{itemize}



\chapter{Technical details and API description}
    \label{ch:faq}

    \section{Usage} Our framework was developed under Linux and therefore it assumes standard Unix tools for development.The source files themselves are supposed to be platform-independent. 

\begin{itemize}
  \item C++11 compliant compiler. Tested compilers are GCC (4.9.3) and CLang (3.7.1).
  \item TinyXML2 library 
  \item Make (optional - for compilation of executables, testing, development)
  \item Dot, gpicview (optional - for graph visualisation)
  \item Standard shell tools (optional - for test scenarios and development) 
  \item Functional bobox environment with aligned memory allocator (optional - for scenario 6)
  \item CPU implementing SSE 4.2  (optional - for test scenarios which generate SSE instructions)
\end{itemize}

The project may be compiled by make:
\begin{code}
> make
\end{code}

By default, make compiles the ctb executable and unit tests and runs basic test scenarios. The following make targets are available:

\begin{description}
  \item[basictest] compiles and runs basic regress tests.
  \item[fulltest] runs full check of the sse instruction table and a test of bobox environment.
  \item[ctb] compiles the ctb executable.
  \item[all] defaults to `ctb' and `basictest'
\end{description}


This library may be used in three ways:
\begin{itemize}
  \item As a separate executable controlled by text commands.
  \item As a header C++11 library. In this case, the \texttt{"ctb.h"} file has to be included in the project. The \texttt{ctb::ctb} class serves as the library frontend. This frontend may be used in two ways, either of them requiring creation of a single object of the ctb instance:
  \begin{itemize}
    \item using text commands
    \item using raw templated calls to commands
  \end{itemize}
\end{itemize}

Basic usage reference of the provided executable may be obtained by:
\begin{code}
> ctb -h
\end{code}


The text command controls of this library may be very easily interfaced by raw C commands, which should allow direct language-independent integration.

The simplest way to invoke ctb is by:
\begin{code}
> ctb -f programme
\end{code}

Here \texttt{programme} is a text file containing a sequence of commands.







    \section{CTB commands} 
The input programme has to conform to the following rules:

\begin{itemize}
  \item Every line represents exactly one whole command unless this line is empty or commented.
  \item Any line can be commented by insertion of \ttt{\#} at its beginning.
  \item Command consists of non-empty space-separated strings. Every sequence of one or more spaces is considered as one delimiter.
  \item The first string in the list identifies the command to be run. The rest of the list is passed as arguments to the command.
\end{itemize}

The available commands are:

\begin{itemize}
  \item \ttt{help} shows a standard UNIX-like help.
  \item \ttt{source <input file>} -- interprets contents of the file as its argument.
  \item Import and export commands. These take the identifier of a loader class as the first argument and a file name as the second argument. The structures are loaded to/exported from the same object of a graph or of an instruction table. The structure is always cleared before load. Note that not all loaders have to implement all import end export functions. 
  \begin{itemize}
    \item \ttt{exportgraph  <loader> <output file>} -- available for the `dot' loader
    \item \ttt{exportinstab <loader> <output file>} -- available for the `csv' loader
    \item \ttt{loadgraph    <loader> <input file>} -- available for the `xml' loader
    \item \ttt{loadinstab   <loader> <input file>} -- available for the `xml' and `csv' loaders
  \end{itemize}
  \item \ttt{generate <output aliasenv> <output file>} -- the \ttt{aliasenv} is invoked with the loaded generator and instruction table. The produced string is written into the \ttt{output file}.
  \item \ttt{testgraph} -- testgraph is a custom command which acts as a special version of loadgraph. This loader function constructs a graph from the loaded instructions, trying to utilise all available operations.  
  \item \ttt{adddebug  [depth=1 [<list of vertex ids>]]} -- This command hooks a debug node\footnote{These are special output nodes which are supposed to output values directly. These are not subject to width conversions.} to every vertex specified in the space separated list of vertex ids as well as to predecessors of these nodes up to \ttt{depth}. Debug nodes are automatically selected from operations of type \emph{debug}. If the list is empty, debug nodes are hooked to all vertices. 
  \item \ttt{transform  <transformer>} -- the graph transformed identified by the first argument is invoked. This transformer is free to edit the loaded graph as it pleases.
  \item \ttt{visualise} -- writes the currently loaded graph into a temporary file and invokes the dot (graphviz) with gpicview to show the graph to the user. We use this mechanism for error reporting as well.
  \item \ttt{preprocess <aliasenv> <filein> <fileout> } -- processes the \ttt{filein} by the aliasenv identified by \ttt{aliasenv} and writes the results to \ttt{fileout}. The \ttt{<aliasenv>} string may be found in a list provided by the \ttt{ctb -h} help.
\end{itemize}


    \section{File structure} \subsubsection{Framework source codes}
\label{sec:directories}
\begin{description}
  \item \ttt{ctb/aliasenv\_*.h} files contain definitions of alias environments.
  \item \ttt{ctb/cartesian\_multiplier.h} contains the cartesian multiplier.
  \item \ttt{ctb/cf\_transform.h} contains graph transformations which perform cycle removal, node expansion and buffer size allocation.
  \item \ttt{ctb/context.h} contains definition of a sample structure interfacing the `cf' aliasenv output.
  \item \ttt{ctb/\string{conversions.h,datatypes.h,defines.h,errorhandling.h\string}} contain auxiliary structures.
  \item \ttt{ctb/\string{graph.h,graph\_factor.h\string}} contain implementation of graph structures.
  \item \ttt{ctb/generator.h} contains implementation of the generator.
  \item \ttt{ctb/instructions.h} contains implementation of structures storing instruction tables.
  \item \ttt{ctb/loader*.h} contain definitions of structures for import, export and possibly dynamic creation of graphs and instruction tables. 
  \item \ttt{ctb/multicont.h} contains our implicit containers.
  \item \ttt{ctb/parser.h} contains an evaluator of arithmetic expressions.
  \item \ttt{ctb/\string{proxy.h,ptrglue.h\string}} contain auxiliary structures serving for syntax-related constructs.
  \item \ttt{ctb/split.h} contains string-processing utilities.
  \item \ttt{ctb/\string{taghandler.h,tagmaster.h\string}} contain definitions of tag-handling structures.
  \item \ttt{ctb/writer.h} contains definitions of the Writer class.
\end{description}

\subsubsection{Other framework files}

\begin{description}
  \item \ttt{ctb/main.cpp} constructs the ctb executable.
  \item \ttt{ctb/test.cpp} constructs an executable containing all unit tests.
  \item \ttt{ctb/Makefile} contains (machine) description of the building process. 
  \item \ttt{ctb/\string{ctb,test\string}} are the produced executable files.
  \item \ttt{ctb/control-flow-notes/} contains the first draft of partition-handling code.
  \item \ttt{ctb/sse\_cf\_macros/buffer.h} contains a set of C macros implementing a SSE in-register buffer. This set is not fully functional.
  \item \ttt{ctb/sse\_cf\_macros/pp\_macros.h} contains macro support for the buffer. These are taken from a third-party repository \cite{cloak}.
  \item \ttt{ctb/sse\_set/src*.csv} contain definitions of SSE instructions. These are split into multiple files to prevent duplicities in different load/store sets.
  \item \ttt{ctb/sse\_set/C\_table*} are tables concatenated from the partial source tables. There is a deprecated version used by the `bobox' environment and a refined version which expects specific handling provided by the `cf' environment.
  \item \ttt{ctb/templates/*} are template files used by some alias environments.
\end{description}

\subsubsection{Test related files}

\begin{description}
  \item \ttt{ctb/tests\_table} contains brief description of test scenarios.
  \item \ttt{ctb/test*/} directories contain test scenarios.
  \item \ttt{ctb/test*/Makefile} are used to execute tests. This may be done by invocation of make.
  \item \ttt{ctb/test*/*.expected} are files which are diff-checked against actual results.
  \item \ttt{ctb/test*/graph.xml} contain description of input graphs.
  \item \ttt{ctb/test*/instab.\string{xml,csv\string}} contain description instruction sets for scenarios. 
  \item \ttt{ctb/test*/\string{program,program\_visual\string}} contain task description for ctb. The visual version usually contains added visualisation commands and arguments. The visual version may be invoked by \ttt{make visual}.
  \item \ttt{ctb/test*/check\_output.sh} contains a sequence of commands which check that results obtained from the tests are correct. (If not present, then this is done by Makefile).
  \item \ttt{ctb/test*/astable.sh} is an awk script which formats streamed tabular information into a tabular view. 
  \item \ttt{ctb/test*/run.sh} present in tests which encompass control flow invokes the test executable altogether with the \ttt{astable} script. This results in a nice visualisation of content of buffers between any two iterations of the crawler algorithm.
  \item \ttt{ctb/test*/macros.h} contain test scenarios in case of scenarios which encompass control flow (12, 13, 16, 17).
  \item \ttt{ctb/test*/*.\string{h,cpp\string}} also contain various test-related functionality such as programmed result-checking, macros for dumping of register content, etc..
\end{description}

\subsubsection{Thesis files}
\begin{description}
  \item \ttt{ctb/thesis/*} are tex template files provided by Martin Mareš. Some of these files have been edited.
  \item \ttt{ctb/thesis/contents/*} are content files.
  \item \ttt{ctb/thesis/graphs/*.dot} contain graph figures of this thesis. These graphs are written in the dot format, i.e., the same format we use for visualisation of graphs in ctb.
  \item \ttt{ctb/thesis/graphs/*.sh} are scripts which preprocess and postprocess the dot files. The primary purpose of postprocessing is making graphs more compact. 
  \item \ttt{ctb/thesis/graphs/fixed/*.tex} are graphs whose tex code needed to be edited by hand. These files are used instead of those of same name generated by dot.
  \item \ttt{ctb/thesis/\string{Make,contents/\string{Make,graphs/Make\string}\string}} take care of building process of this thesis. 
\end{description}




    \section{Text preprocessor} This section explains how the writer (also \emph{the preprocessor}) works. This may be of interest since all text is handled using the writer class. Also the instruction tables specified in csv are preprocessed by the writer.

As we have already indicated, the writer class is basically a string container which allows appending of expressions via a printf like function. Unlike printf, the print method uses shell-like dollar expressions. This allows reuse of arguments and does not restrict the order of arguments. A very simple example follows.

\mybeginfig
\begin{code}
writer w;
w.print("There were $1 $2", 7, "dogs");
w.write_str()
\end{code}
yields
\begin{code}
"There were 7 dogs"
\end{code}
\myendfig{printhello}{A simple example of usage of writer's print method}

There is a number of expansions which take place, most of them recursively when the string is being writen into a writer object.

\section*{Expansions}

The order of expansions is the following:
\begin{enumerate}
\item Cartesian expansion. This expansion is performed only once.
\item All other types of expansions. The \texttt{\$\string{\string}} and \texttt{\$[]} expansions are carried out by a new writer with the same alias environment.
\end{enumerate}

\subsubsection{Argument expansion}

\texttt{\$<n>} (e.g., \texttt{"\$2"}) is resolved to the $n$th argument. This argument is converted to a string and printed using a recursive call which again performs all expansions except for the cartesian expansion.

\mybeginfig
\begin{code}
writer().print("a[$1] = a[$1] $ $2", i, j).write_str()
\end{code}
yields
\begin{code}
"a[i] = a[i] $ j"
\end{code}
\myendfig{printhello}{Another simple example of usage of writer's print method}

\subsubsection{Alias expansion}

Patterns resolved by the alias expansions take the form \texttt{\$<[a-zA-Z]+>}. The identifier is resolved using an alias environment. We should remind the reader that environments are static classes, which are passed to the writer by means of a template parameter.

%aliasenv to resolve the supplemented name - e.g. print("\$animal is an animal") -> "dog is an animal" if the supplemented aliasenv (M) contains alias animal resolved to dog

\subsubsection{Cartesian expansion}
Cartesian expansion takes place only once at the beginning of evaluation of the print call. First, the format string is scanned for all expressions in the form of \\ \centerline{\texttt{\$\string{ <identifier> -> <value1>, <value2> ...\string}} or} \\ \centerline{\texttt{\$\string{ <identifier1>, <identifier2> ... -> <value1>, <value2> ...\string}}} \\ This expression is not allowed to contain nested curly braces. Every such expression is transformed into a list of tuples and replaced by \texttt{\$<identifier1>}.  Then, a cartesian multiplier is constructed from all these lists. Finally, the writer iterates over the product. Per every iteration a new auxiliary environment (which inherits the original environment) is constructed, filled by the currently observed result of the cartesian product and used to print the entire format pattern.

This means that per every line of input, multiple lines of output are emitted, featuring all combinations of variables specified in the lists.

The left side of the expression, i.e., the identifiers, declares a $tuple$ of aliases (constant variables). The first $|tuple|$ values from the right side are then assigned to this tuple and printed. This is repeated while there are any values remaining on the right side.

\mybeginfig
\begin{code}
writer().print(
  "${ color -> red, blue } $animal; nice ${ animal -> dog,cat };"
).write_str()
\end{code}
yields
\begin{code}
"red dog; nice dog\n"
"red cat; nice cat\n"
"blue dog; nice dog\n"
"blue cat; nice cat\n"
\end{code}
\myendfig{cartprint0}{Example of the cartesian expansion}

\mybeginfig
\begin{code}
writer().print(
  "${ color,animal ->blue,dog,red,cat}ish $animal"
).write_str()
\end{code}
yields
\begin{code}
"blueish dog\n"
"redish cat\n"
\end{code}
\myendfig{cartprint0}{Example of tuples in the cartesian expansion}

 %Note that here we do NOT allow nested {}s

\subsubsection{Recursive expansion}
Expressions of the form form \texttt{"\$\string{ whatever that does not contain `->' \string}"} are evaluated in two steps:
\begin{enumerate}
  \item The content of braces is expanded.
  \item The result is evaluated as a standard alias. That means that the result is evaluated by some alias environment and then printed into the current writer. Of course, the final print is again performed with all expansions.
\end{enumerate}

This expansion allows correctly paired nested curly braces.

\subsubsection{Arithmetic expression expansion}

Expressions of the form \texttt{\$[ <expression expression> ]} denote arithmetic expressions. First the content is evaluated by a new writer object. Then, the resulting arithmetic expression is evaluated by the Parser class.

Supported operators are: $(, ), -, +, *, /, \%$


\subsection*{Import and export modes}

\begin{rem}
  This subsection describes more advanced features of the preprocessing environment. Knowledge of these is not required for basic use of ctb.
\end{rem}

If the writer is used to import or export expressions which are to be later processed by another writer object, careful handling of dollar expressions is crucial. For this purpose, we introduce four import/export modes. The writer takes two modes as template parameters. The first mode specifies how the writer should behave when writing data into its private storage. The second mode describes how it should behave when writing them out.

The lifecycle of an expression is the following:
\begin{enumerate}
  \item Print is called.
  \item Cartesian expansion is performed, resulting in a sequence of calls into an internal print method.
  \item Expression is recursively evaluated - single dollar expressions may get expanded or not, depending on the first import/export parameter. During this step, double (and more) dollar expressions are handled according to the first import/export parameter. 
  \item Expression is appended to the internal data container. Also, the expression is broken into lines using the member language type of the root alias environment.
  \item Expression remains in the internal storage until data output is requested.
  \item Content of the writer is outputted. During this phase, all lines are indented and dollar expressions are \emph{again} processed in order to either double, ignore or reduce the multiplicity of dollar signs. This time, the second mode argument is used.
\end{enumerate}

We show the semantics of import/export modes in the following table:

\mybeginfig
\begin{center}
\begin{tabular}{c|c|c|c|c}
                  & Ignore & Eat          & Let          & Expand    \\
  \hline
  \$ expression   & \$     & is expanded  & is expanded  & \$\$      \\
  \hline
  \$\$ expression & \$\$   & \$           & \$\$         & \$\$\$\$  \\
\end{tabular}
\end{center}
\myendfig{dollarexpansions}{Semantics of import/export modes of the writer class.}

  Hence, typical combination will be \emph{Let+Eat} for import operations and \emph{Let+Expand} or \emph{Ignore+Expand} for export. 

\subsection*{Formatting and language}

Formatting is driven by some specialisation of the language structure which was already mentioned. Language structure is supposed to provide the following two methods. These methods are called on every character of the printed string.
\begin{itemize}
  \item \texttt{void shouldindent(\\  string line, int\& outindent, \\  int\& indent, int\& nobreak\\)}\footnote{The ampersand denotes references. In this context, it means that ampersanded arguments act as return values. Besides that, in the actual implementation, we, of course, use constant references to pass the potentially long strings.}
	\begin{description}
	  \item[line] is the line to be indented. 
	  \item[outindent] is the actual indentation of the line.
	  \item[indent] is a persistent indent context.
	  \item[nobreak] allows suppressing of the next \emph{nobreak} line breaks.
	\end{description}
  \item \texttt{void shouldbreak(\\  int pos, string line, \\  bool\& break\_before, bool\& break\_after\\)}
	\begin{description}
	  \item[pos and line] provide an exact position and its context.
	  \item[break\_before and break\_after] are returns which indicate whether the string should be broken before of after the position identified by the \emph{pos} argument in \emph{line}.
	\end{description}
\end{itemize}





    \section{Format specifications}

      \label{sec:formats}

Our basic input method is input in form of xml files. Besides xml file, we implement also a plain text input interface for instruction tables since this flat files turned out to be much easier to create. 

Regarding output methods, our framework creates neither graphs nor instruction tables, so we do not have any motivation to provide any implicit output interface. All output interfaces we do implement we implement for sake of development convenience.

\subsection{XML}

We use standard XML format, as parsed by the TinyXML2 library. Any field of any structure may be represented either by attribute or by a separate child element. Of course, fields which are present in multiple instances have to be present as child elements.

The structure is described by the following tree. The attributes which are not supposed to be present in multiple instances are specified as the lists following their parent elements.


Graph:
\dirtree{% 
.1 graph\_list.
.2 graph.
.3 vertex - params, vid, opid.
.3 edge - from, to, from\_pos, to\_pos.
}

\pagebreak

Instruction table:
\dirtree{%
.1 root.
.2 instruction\_list.
.3 operation - input, output, debug, opid, out\_type, in\_type.
.4 instruction - width\_in, width\_out, code, note, tags, \\ rating.
.5 code\_custom .
.2 expansion\_list.
.3 expansion - opid, name, transformer\_name, arguments, note.
.2 type\_list.
.3 type - typeid.
.4 type\_version - width, bitwidth, code, note.
.4 type\_conversion  - width\_in, width\_out, code1, code2, \\ code\_custom, code\_generic, note, tags, rating.
}

The following structures are expected to be present at most once: \\ \centerline{\ttt{root, instruction\_list, expansion\_list, type\_list, graph\_list, graph}}

The exact meaning of fields will be given in one of the following subsections.

\subsection{CSV}

Our format has the following general properties:
\begin{itemize}
  \item Every line stands for one record unless this line is empty, commented or joined with the following line.
  \item Every record consists of tab-separated fields.
  \item A line can be commented by \ttt{\#} put at the very beginning of the line.
  \item Any line that has the \ttt{\\} at its end is joined with the following line.
  \item Our use of this format has the following semantics:
  \begin{itemize}
	\item Fields which represent lists are comma separated.
	\item This file is preprocessed on load by a text preprocessor\footnote{An instance of the Writer class parametrized by some suitable parameters.} when the file is loaded.
  \end{itemize}
\end{itemize}

Although we use the term CSV, the reality is that our flat plain-text format has almost nothing in common with the RFC 4180\footnote{RFC stands for Request For Comment and is a standard means of publishing internet standards.}.

The reason why we have decided to add a flat representation of instruction tables is that flat plain-text formats are much easier to process by general text-processing tools. Specifically, our problem is that a single instruction table may contain thousands of instructions, which typically share some patterns in some fields and mainly the fact that someone has to write them. For this reason, we have decided to create a flat format which would be preprocessed by our own custom preprocessor. Specifically the cartesian expansion was designed specially for the purpose of simple description of instructions.

This approach, of course, presents one difficulty --- we wish to encode a tree structure into a flat file. We flatten the structure by using one line per every leaf of the tree while specifying significant information of the entire path. This way, some information may be present duplicitly. In reality, we often wish to build instruction tables bottom-up, using a single scheme per multiple leafs, therefore generating multiple instructions of (typically) \emph{different} operations. 

The motivation presented in the previous paragraph may need some example. Imagine that we want to generate multiple vectorized versions of the bitwise \emph{or} operation. We will most likely use the same instruction pattern (the same intrinsic function) and just specify it for all integer types which we have defined (signed and unsigned integers of varying bit lengths). This way, we generate one instruction per every integer \emph{or} operation. Note that these are different operations since they operate on different types.

We distinguish multiple different types of lines. We do so by the second column, which contains type of record. The first column serves for user-defined note, which serves for user's orientation in the file, development notes, debug...

We show the structure by the following tree. In this tree, every path between the root node and any leaf list specifies one possible structure of a line. The fields found on this path must be present in their order.

\dirtree{%
.1 note.
.2 type=instruction.
.3 \text{\emph{(operation definition)}} out\_type, in\_types, opid, flags\\ \text{\emph{(instruction definition)}} width\_in, width\_out, code, tags, \\rating, code\_custom, code\_custom, code\_custom....
.2 type=expansion.
.3 \text{\emph{(operation definition)}} out\_types, in\_types , opid \\ \text{\emph{(expansion definition)}} name, transformer\_name, args.
.2 type=type\_version.
.3 \text{\emph{(logical type definition)}} type\_id, bitwidth \\ \text{\emph{(version definition)}} width, code.
.2 type=type\_conversion.
.3 \text{\emph{(logical type definition)}} type\_id \\ \text{\emph{(conversion definition)}} width\_in, width\_out, code\_1, \\ code\_2, code\_custom, code\_generic, tags, rating.
}

\subsection{Detailed field semantics}

\subsubsection{Common}
\begin{description}
  \item\ttt{note} is an arbitrary user-defined string. This is not used by the framework at all, but is preserved for potential debug purposes or for export.
  \item\ttt{tags} is a comma separated list of user-defined strings. These serve for restriction of the instruction set. These flags may be restricted through the command line interface or directly from an alias environment.
  \item\ttt{rating} is an integer which allows definition of precedence. Always the instruction with highest rating is used if there are more feasible instructions for the required width.
\end{description}

\subsubsection{Operations and instructions}
  \begin{description}
  \item \ttt{flags} (csv field) -- comma separeted list consisting of flags from the following list. These identify special types of operations. These nodes may be treated specially by graph transformations and also affect how the default \ttt{code} field is used.
  \begin{description}
	\item \ttt{input} -- operation which acts as a data source. 
	\item \ttt{output} -- operation which acts as a data sink.
	\item \ttt{debug} -- operation which acts as a special data sink for the \ttt{adddebug} command.
  \end{description}
  \item \ttt{input, output, debug} (xml boolean flags) -- correspond to the previous field of csv.
  \item\ttt{opid} -- identifier operation.
  \item\ttt{out\_type} -- \ttt{type\_id} of the output type. The types of operations are strictly checked.
  \item\ttt{in\_types} -- list of \ttt{type\_id}s. The types of operations are strictly checked.
\item\ttt{width\_in} -- input width of an instruction (i.e. number of independent data rows its instruction operates on).
  \item\ttt{width\_out} -- output width of an instruction. Operations are not currently allowed to perform implicit width conversions, so this field is always equal to \ttt{width\_in}.
  \item\ttt{code} -- an instruction pattern for the instruction. This field is used with respect to the type of node, which means that it is substitued into one of three predefined patterns. For instance \ttt{code} may be \ttt{"\$arg1 + \$arg2"}. Note that dollar signs must be doubled in case of the csv format since the csv loader expands single-dollar expressions.
  \item\ttt{code\_custom} -- taking the format of \ttt{"\textless identifier\textgreater:} \ttt{\textless instruction pattern \textgreater"}. The pattern will be printed into a layer of writer which is identified by \ttt{<identifier>} --- for this reason it may be present multiple times. Custom code does not use a predefined code template, so full code must be present. For example \ttt{"\$type \$name = \$arg1 + \$arg2;"}.
\end{description}

\subsubsection{Types, type versions}
\begin{description}
        \item\ttt{type\_id} -- an identifier of a type.
        \item\ttt{width} -- width operations which operate on this type version. Or length of this vector.
        \item\ttt{bitwidth} -- hardware bit length of this type per one record. This allows nontrivial extension-specific computations in instruction patterns. This may be 1 or 8 for a boolean type independently of the \ttt{width} argument.
        \item\ttt{code} instruction pattern for this type. For instance \ttt{"int"} for an integer.
\end{description}


\subsubsection{Type conversions}
\begin{description}
          \item\ttt{width\_in} -- width of input vectors.
          \item\ttt{width\_out} -- width of output vectors.
          \item\ttt{code1, code2} if $width\_in/width\_out = 2$, this two field act as selection expressions of the first and second half of the input vector. If $width\_out/width\_in = 2$, then the second field is ignored and the first one acts as a merge expression merging two vectors into one. We show this in example:

\mybeginfig
\begin{code}
width_in=1
width_out=2
code1="$arg0 | ($arg1 << 16);"
\end{code}
\myendfig{inspast0tex}{Width conversion corresponding to line 3 of figure \ref{fig:vectorizedfragment}}

\mybeginfig
\begin{code}
width_in=2
width_out=1
code1="$arg1 & 0xFFFF"
code2="($arg1 & 0xFFFF0000) >> 16"
\end{code}
\myendfig{inspast0tex}{Width conversion corresponding to lines 5 and 6 of figure \ref{fig:vectorizedfragment}}
          \item\ttt{code\_custom} this is a nonmanaged version of the field. This means that all names, types and conversions have to be present in this single field, if this field is to be used on its own.
          \item\ttt{code\_generic} allows iterated select or insert operation. This means that we may use this field if we can describe all expressions by the same pattern. For instance, if we want to join four elements into a single vector, we may declare the target variable by the \ttt{code\_custom} field and then write a generic select which will modify the variable once per every call.  For this reason, the \ttt{code\_generic} is guaranteed to be generated after \ttt{code\_custom}.

\mybeginfig
\begin{code}
code_custom="$type $name[4];"
code_generic="$name[$vindex] = $arg;"
\end{code}
\myendfig{instex}{Example of use of the \ttt{code\_generic} field showing join of four values into one vector.}

\end{description}

\subsubsection{Graph fields}
\begin{description}
\item\ttt{params} -- comma separated list of assigment such as \ttt{"var1=dog,var2=cat"}. Hereby specified values are available during code generation. In other words these values may be used in instruction patterns.
\item\ttt{vid} -- identifier of a vertex.
\item\ttt{opid} -- operation identifier.
\item\ttt{from} -- vertex id of an edge.
\item\ttt{to} -- vertex id of an edge.
\item\ttt{to\_pos} -- the $to$ annotation identifying arguments of operations.
\item\ttt{from\_pos} -- the $from$ annotation.
\end{description}

\subsection{Instruction pattern fields semantics}

This subsection describes aliases provided by the generator, i.e., `variables' which may be used in instruction patterns. Of course, any other aliases may be defined in custom alias environments and used in instruction patterns.


\subsubsection{Generator-specific aliases}

\begin{description}
\item\ttt{type} -- pattern of the type of the result for the correct width.
\item\ttt{name} -- name generated for the variable.
\item\ttt{name1, name2...} -- names generated for the variables in the \ttt{code\_custom} field of split operation.
\item\ttt{classid} -- index of the current partition. These are in topological order with respect to the graph structure.
\item\ttt{basename}  -- name of the first variable in vector. 
\item\ttt{operation} -- the instruction pattern which is being processed.
\item\ttt{arg1, arg2, arg3...} -- arguments of the operation in question. These are names memoized from some previous iteration of the generator, generated either from an instruction or from a width conversion.
\item\ttt{iindex} -- input index -- index denoting position of the first data row of the vector in the currently processed pack of data. 

  For instance, for standard instruction of width 4 and pack size 16 the pattern will be expanded four times with values  0, 4, 8 and 12.

\item\ttt{oindex} -- output index -- output analogy of the \ttt{iindex}. These two values will differ in width conversions.


\item\ttt{vindex} -- vector index -- index denoting position of the currently processed subvector in a vector processed during width conversion. In other words, this is always the difference between \ttt{iindex} and \ttt{oindex}.

This value is meant for use with generic joins and splits. For instance, assume width conversion from 8 to 2 at pack size 16. Then the \ttt{code\_generic} pattern will be expanded with values 0, 2, 4, 6, 0, 2, 4, 6. The \ttt{iindex}es will be 0, 0, 0, 0, 8, 8, 8, 8 while \ttt{oindex}es will be 0, 2, 4, 6, 8, 10, 12, 14.
\end{description}


\subsubsection{Other default aliases}

The following patterns are specified as constants in the aliasenv of the generator as default C-specific patterns. Since these are language specific, one may need to override them. This may be done by simple specification of these in your custom alias environment. Those are exactly those patterns which are chosen based on the \ttt{flags} field and affect the \ttt{code} field.

\begin{description}
\item \ttt{"\$innercode" -> "\$type \$name = (\$type)(\$operation);"}
\item \ttt{"\$declcode" -> "\$type \$name;"}
\item \ttt{"\$outputcode" -> "\$operation /*\$name*/;"}
\item \ttt{"\$inputcode" -> "\$type \$name = (\$type)(\$operation);"}
\item \ttt{"\$conversioncode" -> "\$type \$name = (\$type)(\$operation);"}
\end{description}

We also propose use of the following three aliases which lay entirely in user space.

\begin{description}
  \item \ttt{ioindex} -- this field is meant to identify input or output data for the purpose of effects. We propose this index to be set as a parameter in vertices of input and output operations.
  \item \ttt{input}, \ttt{output} -- we propose these fields to be constant patterns set in an alias environment to expressions representing actual input or output values consistent with semantics of the input and output instructions. These patterns would then be used in instruction tables. This proposal ensures that instruction tables are not polluted by environment-specific side effects. 

  For instance, \ttt{"\$input" -> "input\_\$ioindex[j]"}. This pattern assumes that the context filled in by the alias environment responsible for composition of output prepares data into arrays and then iterates over them in a loop with index j. 
\end{description}




    \section{API description} Note that in this section, we will ommit all insignificant information such as constant reference markings or template definitions.

\subsection{Loader API}

Loaders are required to implement the following five methods:

\begin{code}
void load_graph(G& graph, istream&) ;
void load_instab(IT& instab, istream&) ;
void export_graph(G& instab, ostream&) ;
void export_instab(IT& instab, ostream&) ;
static string get_name();
\end{code}

\begin{description}
\item \texttt{G, IT} are templated instruction table and generator types.
\item \texttt{istream, ostream} are C++ implementation of input and output streams. This way, direct interfacing with another programme is possible.
\item \texttt{get\_name} is required to return a sensible name not colliding with names of other loaders.
\end{description}

We believe that provided the description above, the api is self describing.

When a new loader is implemented, it is supposed to be registered within the ctb object by its templated \texttt{register\_loader} method. This call may be inserted into the \texttt{fill} method of ctb. Or may be inserted into a new instance of the \texttt{main.cpp} file. Registration is not neccessary if the user does not wish to use the text command controls.

\subsection{Aliasenv API}

Fully functional alias environment which acts as an output environment is supposed to provide the following interface:

\begin{code}
typedef <typename> language;
static string alias(string a, bool* s = NULL);
static string get_name();
static writer generate(int packsize,  G& generator, 
  string name, stringlist args);
\end{code}

\begin{description}
\item \texttt{language} is a language structure which is to be used by any writer parametrized by this environment.
\item \texttt{alias} is a method which tries to evaluate the string a. If the second parameter is not null, it is also supposed to report whether the attempt was successfull.
\item \texttt{get\_name} is again any sensible string identifier.
\item \texttt{generate}
  \begin{description}
    \item \texttt{packsize} is the size of the pack for which the body is supposed to be generated, i.e., the $w$ from chapters dedicated to theoretical analysis. The environment is allowed to ignore this parameter and decide its own suitable size.
    \item \texttt{generator} is again a templated generator type. This variable provides direct access to the underlying graph as a simple member variable.
    \item \texttt{name} is a simple identifier meant for identification of the generated code fragment. E.g., the alias environment may output its results as member methods of a class whose name contains this identifier.
    \item \texttt{args} is a list which may contain other arguments.
  \end{description}
\end{description}

If an environment is constructed to act just as an alias processor, only the \texttt{alias} method and the \texttt{language} type are obligatory. If an environment is not supposed to be used as the top-level parameter of writer, the \texttt{language} field may be omitted as well.

As other structures, enviroments which provide the \texttt{generator} method are supposed to be registered within ctb by the \texttt{register\_aliasenv} method.

\subsection{Generator API}

Generator provides the following methods for the purpose of code generation. 

\begin{code}
void generate(int packsize, 
    imp_contB<W>& w, imp_contB<output_options>& options); 
void generate_partition(int partition, int packsize, 
    imp_contB<W>& w, imp_contB<output_options>& options);
imp_contB<output_options> option_struct();
int partition_count();
\end{code}

\begin{description}
\item \texttt{w} is the writer which is supposed to receive the output. This type may be constructed as follows:
\begin{code}
impl_contB<writer<example_aliasenv>> w;
\end{code}
 The results in nondefault layers are to be retrieved as follows:
\begin{code}
w["empty_output_indicators"].write_str();
\end{code}
\item \texttt{options} specify special behaviour for specific layers. Instance of this structure may be either manually constructed or obtained by the \texttt{option\_struct} method. Setting options works exactly as in case of the writer. For instance:
\begin{code}
auto o = generator.option_struct();
o["empty_output_indicators"].unique = true;
\end{code}
Currently available options are:
\begin{description}
\item \texttt{once} causes generator to write an instruction into such layer only once per vertex.
\item \texttt{global\_once} causes generator to write the value into such layer only if the corresponding layer of writer does not contain anything yet.
\end{description}
\end{description}

For the purpose of graph transforms, the generator provides the following methods. These methods interface the corresponding methods of the graph member, sometimes handling generation-related data specially. Advanced functionality, such as custom graph crawling is provided by the graph itself.
\begin{code}
node* get_vert(J id);
node* add_vert(v_id v, op_id op, map<string,string> params) ;
void add_edge(v_id from, v_id to, int to_pos, int from_pos) ;
void rm_vert(J v);
void rm_edge(edge* e);
void connect_as(J v, K as, bool inputs = true, bool outputs = true);
void foreach(function<void(node*)> f);
\end{code}

\begin{description}
  \item \texttt{J, K} are automatically resolved types identifying vertices. These may be pointers or references to the actual vertex structure or just vertex ids.
  \item \texttt{v\_id, to\_pos, from\_pos, params} correspond to the described field semantics.
  \item \texttt{get\_vert, add\_vert, rm\_vert, add\_edge, rm\_edge} apparently allow altering of the graph structure.
  \item \texttt{connect\_as} adds new edges to the vertex \texttt{v} imitating the way the \texttt{as} node is connected.
  \item \texttt{foreach} calls the function \texttt{f} once per every vertex in topological order. This function may be nested.
\end{description}

The retrieved nodes then provide the following api:
\begin{code}
T data; I id;
int classid, colourmark; 
imp_cont<vector<edge*>> in, out;
edge* in_at,out_at(int i, bool check_uniqueness = true);
\end{code}

\begin{description}
  \item \texttt{id} contains vertex id.
  \item \texttt{classid} is id of the component of connectedness if the factor graph is up-to-date. Make sure you have called the \texttt{graph.update\_factor()} method. This is not done automatically since factor graph and this id is apparently not stable with respect to factorisation.
  \item \texttt{colourmark} is and auxiliart identifier which may be used by the user for marking of vertices during graph algorithms.
  \item \texttt{data} contain the data type specified by the \texttt{T} template parameter. The generator's data type contains a pointer to the corresponding operation (\texttt{op}), an operation id (\texttt{opid}) and parameter access methods (\texttt{get\_param(string)} and \texttt{set\_param(string,string)}). 
  \item \texttt{in, out} are containers containing pointers to incoming and outgoing edges. Layer 0 may be iterated by simple \texttt{for(e : in)}, other layers by \texttt{for(e : in.get\_layer(1))}.
  \item \texttt{in\_at, out\_at} provide direct access to the \texttt{i}th edge.
\end{description}

\subsection{Instruction table import API}

Instruction tables may be constructed via the following api:

\begin{code}
operation& add_operation(opid_t op_id, tid_t out_type, 
  vector<tid_t>& in_types, flag_t flags) ;
void operation::add_code(int width_in, int width_out, string code, 
  stringlist code_custom, string note, string tags, int rating);

type& add_type(typename T::tid_t t, int bitwidth = 0) ;
void type::add_code_type(int width, string code, string note);
void type::add_code_conversion(int width_in, int width_out, 
  string code1, string code2, string code_custom, 
  string code_generic, string note, string tags, int rating);

void add_expansion(opid_t op_id, string name, string transformer, 
  vector<opid_t> args, string note, vector<tid_t> in_types, 
  vector<tid_t> out_types);
\end{code}

\begin{description}
  \item \texttt{add\_operation, add\_type} look up the requested operation or type. If the object does not exist, it is created. Finally reference to this object is returned, allowing insertion of child elements.
  \item other fields are exactly the fields described by field semantics.
\end{description}

\subsection{Graph transformer API}

Api required from custom transformers is the following:
\begin{code}
 static string get_name();
 void transform(G& generator, stringlist args);
\end{code}

Semantics of these fields are again the same. There are no restrictions on content of \texttt{args} and no restrictions on what the transformer is allowed to do with the generator or with its graph.


\subsection{Custom commands}

Custom commands may be registered within the ctb class by the following method:
\begin{code}
void register_command( string cmd, function<void(stringlist&&)> f,
  string description);
\end{code}

\begin{description}
\item \texttt{cmd} -- the command string.
\item \texttt{f} -- function performing the command.
\item \texttt{description} -- short reference which will by shown in the '-h' option.
\end{description}

  




    \section{Testing} \subsection*{Framework testing}

\subsubsection{Unit testing}
Every unit of ctb contains a static \texttt{self\_test} method. This method contains some scenarios which are supposed to ensure that everything works properly. All these methods are included by the \texttt{test.cpp} file and compiled as a separate executable.

\subsubsection{Test scenarios}
We also provide a number of scenarios which should ensure that the composed whole works as expected. Some of these scenarios test graph transformations, some test alias environments, some test implementation of macros, some test instruction tables...

These tests are usually specified by means of a \texttt{program} file. Results are usually checked agains remembered outputs which were verified by a human being. Since order of some iteration methods of the C++ language is not stable across different compilers, some outputs may differ.

The test directories also serve as examples of usage of ctb.

\subsection*{Instruction table testing}

For the purpose of testing of instruction tables, we provide the \texttt{testgraph} command, which generates a new graph such that:
\begin{itemize}
  \item Multiple instances of every input operation are generated.
  \item An instance of every data-processing\footnote{Non-input-non-output operation.} operation is created. Inputs are connected to these operations so that different arguments always receive input from different inputs.
  \item An output operation is created for every output of a data-processing operation.
\end{itemize}

When writing a new instruction table, we usually wish to ensure two things:
\begin{itemize}
  \item Syntactical correctness of regular instructions.
  \item Semantical correctness of regular instructions.
\end{itemize}

Syntactical correctness may be ensured by means of compilation of code generated from the testing graph. However, experience shows that these graphs are of enormous size, and therefore their compilation may take very long time. As a solution, we provide a special option \texttt{-c}, which makes the compiler output \emph{only} the first instruction per every vertex. For all other instructions, only a dummy variable is declared.

For testing of semantical correctness of instruction tables, we have implemented the `simu' environment. This environment outputs a fragment of code which constructs an input using a pseudorandom generator. When the resulting c file is compiled and run, it repeatedly calls functions realizing the graph by different widths and checks that all results are equal to the width 1 results. If a collision is encountered, the problematic fragment is invoked again with a break-point-like indicator set. This causes that detailed log is generated during the iteration which observed wrong behaviour. This log contains contents of sse registers generated via \emph{debug} nodes.

We do not cover testing of special control-flow-related instructions since these require tailored testing scenarios.


\subsubsection{Scenarios}

Test directories of our framework are numbered. The following scenarios may be found in the project:

\begin{enumerate}
  \item[1] consists of the unit tests and some simple hard-coded generation. The input for this generation is based on sse but is not meant to provide functional output. This results are produced by the `simple' environment.
  \item[2,3] test that loads end exports of tables work properly.
  \item[4,5] are `simu' environment tests. The 4 is a full test of the provided SSE table. The 5 shows the same functionality on a simple graph and unlike the 4th test is part of basic test chain.
  \item[6] provides output in form of bobox boxes. This environment shows full implementation of preloads with respect to bobox semantics of envelope usage.
  \item[7, 8] test cycle removal and node expansion.
  \item[10] tests functionality of provided SSE buffer macros.
  \item[11] shows functionality of non-vectorized control flow on a simple C table with hard-coded input-output operations. This scenario uses the buffer macros with simple, nonvectorized variables.
  \item[12] shows the same functionality with fully vectorized bodies.
  \item[13] shows the same functionality, but uses fully vectorized buffers.
\end{enumerate}

The test 9 was meant to test the C preprocessor implementation of a SSE buffer which may be found in the \texttt{sse\_cf\_macros} directory. If the reader is interested in advanced usage of C preprocessor, he is welcomed to look around. Yet first, we would like to give him one advice:
\\ \\ \centerline{\sl{``Abandon all hope, ye who enter here.''}\cite{TODO}}




    \section{Available environments} The following environments with generation capabilities are present:

\subsubsection{Simple}
`Simple' environment was the first implemented environment. It prepares random data into arrays created according to the graph specification and runs the result for a fixed number of iterations. This environment handles only graphs consisting of regular operations.

\subsubsection{Simu}
`Simu' environment is an improvement of the `simple' environment. This environment was designed for testing of the provided SSE instruction table.

\subsubsection{Bobox}
`Bobox' environment provides functionality similar to the simple environment except for two things. First, its API is API of the bobox environment and second, it is tailored to utilize the capabilities of aligned load and store operations. Besides the generator-related functionality, this environment also generates code which manages the envelope system of the bobox environment. With hindsight, we have to admit that this API layer should rather have been abstracted entirely outside of our framework. 

  \subsubsection{Cf}
  `Cf' environment implements algorithm \ref{alg:crawler}. This environment already abstracts its API and fully uses the SSE instruction set.
    
    The output file consists of four functions enclosed into a class with two templated context types (this is shown in figure \ref{fig:fragmentapi}). These four functions take a number of rows to be processed and a context structure which provides arrays of data, indices and offsets. The required context structure API is shown in figure \ref{fig:contextapi}. 

\mybeginfig
\begin{code}
template<typename CTXIN, typename CTXOUT>
class box_<name>
{
  void process_single    (int, CTXIN&, CTXOUT&);
  void process_aligned   (int, CTXIN&, CTXOUT&);
  void process_unaligned (int, CTXIN&, CTXOUT&);
  void process_shifted   (int, CTXIN&, CTXOUT&);
}
#define IN_INDICES_<name> <values>
#define IN_TYPES_<name> <values>
#define IN_TYPES_IL_<name> <values>
#define OUT_INDICES_<name> <values>
#define OUT_TYPES_<name> <values>
#define OUT_TYPES_IL_<name> <values>
#ifdef RUN_BOX
  RUN_BOX(box_<name>, IN_INDICES_<name>, IN_TYPES<name>...);
#endif

\end{code}
\myendfig{fragmentapi}{API of the resulting code of the `cf' environment.}


\mybeginfig
\begin{code}
struct context
{
  <type1>* data_0; <type2> *data_1; ...
  int index_0;     int index_1;  ...
  int offset_0;    int offset_1; ...
}
\end{code}
\myendfig{contextapi}{API of the context for the `cf' environment.}


    The \ttt{\$input} pattern of this environment is set to \ttt{"ictx.data\_\$ioindex[ictx.index\_\$ioindex+\$iindex]"}. Indices are supposed to be incremented by entire packsize at ends of code of partitions --- this way we do not introduce new dependencies between load instructions. Indices may be simply incremented by adding the \ttt{"\$inputinc"} pattern to the \ttt{"increments"} pattern layer. The alias \ttt{inputinc} is set to \ttt{ictx.offset\_\$ioindex+=\$packsize}. Aliases \ttt{"output"} and \ttt{"outputinc"} are defined analogically.

    The following writer layers are used for the purpose of instruction-code generation:
    \begin{description}
      \item \ttt{global} gets outputted into the toplevel scope of the processing function. Suitable for injection of testing values.
      \item \ttt{default} layer serves for the standard output.
      \item \ttt{preload} layer has the \ttt{once} flag turned on and is directed to the default layer, meaning that this layer is generated before the \ttt{default} layer while the topological ordering of all instructions is preserved.
      \item \ttt{shiftacum} and \ttt{shiftacumpreload} are analogical to the pair of \ttt{default} and \ttt{preload}. The difference is that this layer gets outputted into a different layer which is handled specially. Namely, when the \ttt{"EXPANDSHIFT"} alias is encountered in any output layer, the content of the acumulator is printed into a swich pattern which realizes shifted loads. The resulting switch gets printed into the default layer.
    \end{description}

Also, there are also 7 macro constructs which are designed to enable easy automatic management of the resulting box. Namely these contain:
\begin{description}
  \item \ttt{IN\_INDICES\_<name>} set to coma separated list of input indices.
  \item \ttt{IN\_TYPES\_<name>} set to coma separated list of input types.
  \item \ttt{IN\_TYPES\_IL\_<name>} set to coma separated list of input types and interleaved by dummy types in index order. This means that if this list is zipped with the sequence \ttt{0, 1, 2, 3...} then the types correspond to their \ttt{ioindex}es.
  \item Analogical macros for outputs.
  \item \ttt{RUN\_BOX} expansion if the macro is defined. This macro may be simply included into the resulting code fragment. In that case, the macro may use the defined index and type lists to construct appropriate context types and to pass these into some processing function.Such function is then free to fill the contexts with data and use the generate box as it finds appropriate. 
\end{description}

We also provide a sample context template (\ttt{context.h}). This template provides storage for the required fields, basic methods for offset calculation and also acumulator lists. The lists allow generic processal of these contexts (such as filling of input fields and retrieval of the processed data). In other words, any processing environment (such as bobox) can be simply interfaced by a single file containing a single generic algorithm. There is no need to generate environment-specific interface for every box.






    \section{Implementation of new environments} Implementation of an entirely new environment, e.g., providing support for an entirely new language, a SSE extension and a target distributed environment, may require the following steps:

\begin{enumerate}
  \item Implementation of a new instruction table.
  \item Implementation of a new alias environment. 
  \item Implementation of a new language formatter, unless you are fine with non-formatted output files. In that case feel free to use the \texttt{language\_empty} structure.
  \item If a new control flow strategy or new types of control flow nodes are to be introduced, a new graph transformer may also be neccessary.
\end{enumerate}

The first two things which the user should think about are:

\begin{itemize}
  \item Representation of vectorized data types. All elements need to retain the correct order when being split or merged.
  \item Exact strategy of employment of control flow.
\end{itemize}








\chapter{Implementation of the SSE instruction set}
    \label{ch:sse}

    We will skip any broader introduction of the SSE extension, since deeper knowledge of the actual instruction set is not necessary for the purpose of this text. A very brief and very informal introduction to the SSE extension may be found in \cite{ssetut}. SSE reference documentation is provided by \cite{reference}. A handy reference is the cheat-sheet available at \cite{cheats}. Paper which discusses more advanced aspects of efficient use of the SSE (and similar) extensions is \cite{autotuning}. We provide only the necessary conceptual minimum of the SSE type system.

We will omit the \ttt{\_mm\_} prefixes and other non-substantial information contained in function names.

\section{Data types and conversions} 

\subsubsection{SSE type system}

SSE provides a number of 128 bit registers, which may be directly used by means of the \ttt{\_m128} type. Such register can contain $128/\mf{length}$ values of signed or unsigned integers of lengths from 8 to 64, single-precision floating point numbers and double-precision floating point numbers. Boolean values are supposed to be represented by masks whose bits are either all 1 or 0.

The content of these registers is to be interpreted as a right-to-left list of values. Note that semantics of shifting operations make this point important. The reader may ask why is it to be this way, and rightly so --- the reason is the clash of our customs of writing lists left-to-right while writing numbers right-to-left. Consistent, increasing order of elements is often desired for the purpose of width conversions and casts.

\subsubsection{Our type system}

As we have already mentioned, we provide the following data types:
\begin{itemize}
\item Signed and unsigned integers of widths 8, 16, 32, 64.
\item Single-precision and double-precision floating point numbers.
\item Byte bools, which are suitable for use with control flow.
\item Bit Boolean vectors, which are suitable for high throughput computations without control flow.
\end{itemize}

We use the simplest and most obvious way of arrangement of data in SSE registers when a register is to be filled entirely. However, when two vectors of different widths of the stored data types are to be processed by one operation, we need to select only some data from the vector of the narrower data type width. Broadly speaking, we always use positions in their intuitive order which are multiples of $128/n$ where $n$ stands for the number of data elements we want to store into a single register. 

Consider the following example of a 16 bit integer width conversion from width 8 to 4 (numbers represent order of elements).


\mybeginfig
%\begin{center}
\begin{longtable}{ c c c }
  {
\begin{tabular}{|c|c|c|c|c|c|c|c|}
\hline
8 & 7 & 6 & 5 & 4 & 3 & 2 & 1 \tabularnewline
\hline
\end{tabular}
}
&
$\arrow$
&
  {
\begin{tabular}{|c|c|c|c|c|c|c|c|}
\hline
\ \ & 8 & \ \ & 7 & \ \ & 6 & \ \ & 5 \tabularnewline
\hline
\ \ & 4 & \ \ & 3 & \ \ & 2 & \ \ & 1 \tabularnewline
\hline
\end{tabular}
}
\end{longtable}
%\end{center}
\myendfig{sse_conversion}{Width conversion realised on SSE registers.}

This way, elements stay always aligned with other elements from the same data row. Such vector may enter a computation with another vector consisting of four 16-bit integers. But it may as well enter a computation with a vector of four 32-bit integers. This property is especially handy with bools if we edit Boolean splits to copy the value into both fields. Visually this means performing the previous split as follows:

\mybeginfig
%\begin{center}
\begin{longtable}{ c c c }
\begin{tabular}{|c|c|c|c|c|c|c|c|}
\hline
8 & 7 & 6 & 5 & 4 & 3 & 2 & 1 \tabularnewline
\hline
\end{tabular}
&
$\arrow$
&
\begin{tabular}{|c|c|c|c|c|c|c|c|}
\hline
8 & 8 & 7 & 7 & 6 & 6 & 5 & 5 \tabularnewline
\hline
4  & 4 & 3 & 3 & 2 & 2 & 1 & 1 \tabularnewline
\hline
\end{tabular}
\end{longtable}
%\end{center}
\myendfig{sse_boolean_conversion}{Width conversion suitable for bools.}

Notice that to perform this split, a nontrivial data reordering was necessary. A simpler approach would simply select every second element by an \emph{and} mask and then put it into a new vector. Showing this by a pseudo-vector C operation, we would do:

\mybeginfig
\begin{code}
int result1 = (input & 0x00FF00FF);
int result2 = (input & 0xFF00FF00) >> 8; 
\end{code}
%\begin{center}
\begin{longtable}{ c c c }
\begin{tabular}{|c|c|c|c|}
\hline
 4 & 3 & 2 & 1 \tabularnewline
\hline
\end{tabular}
&
$\arrow$
&
\begin{tabular}{|c|c|c|c|c|c|c|c|}
\hline
 \ \ & 4 & \ \ & 2 \tabularnewline
\hline
 \ \ & 3 & \ \ & 1 \tabularnewline
\hline
\end{tabular}
\end{longtable}
%\end{center}
\myendfig{sse_bad_conversion}{Example of a possibly malfunctioning width conversion.}

The results are still aligned. The problem with this operation comes when two vectors of different base width come into contact. If we wanted to sum the vectors from Figure \ref{fig:sse_bad_conversion} with a vector of 16-bit integers (the original width of \ttt{input} was 8-bit), we would get an element-wise sum of the vector $((4,2),(3,1))$ with the vector $((4,3),(2,1))$. Of course, this problem may be further addressed by proper setting of a fixed ordering per every data width. The technical difficulty is that all load, store, split and merge operations would have to shuffle data in a rather nontrivial fashion. 

%https://software.intel.com/sites/default/files/a6/22/18072-347603.pdf

Width conversions (as described above) may be implemented by means of the set of \ttt{shuffle} instructions provided by SSE. We provide an excerpt dedicated to shuffle macros from the SSE documentation \cite{reference} to portray its semantics: 

\myloosequote{ The Streaming SIMD Extensions (SSE) provide a macro function to help create constants that describe shuffle operations. The macro takes four small integers (in the range of 0 to 3) and combines them into an 8-bit immediate value used by the SHUFPS instruction. You can view the four integers as selectors for choosing which two words from the first input operand and which two words from the second are to be put into the result word. }{reference}





    \section{Considering load and store problems}  
The SSE extension provides fast aligned load and store operations. These operations load vectors of values from continuous memory areas which are required to begin at multiples of 16 bytes (128 bits). 

The extension also provides the \ttt{alignr} operation. This operation takes two SSE registers and an \emph{offset}, concatenates the two registers and selects a new (continuous) SSE register shifted by \emph{offset}. 

The \ttt{alignr} opperation allows implementation of an efficient non-aligned load. This works by preloading the first aligned area into an auxiliary variable and then chaining a sequence of load operations through the \ttt{alignr} operation. 

\mybeginfig
\begin{code}
_m128 aux1, aux2;
_m128 aux1 = load(input[j-offset]);

aux2 = load(input[j-offset+1*packsize]);
_m128 a1 = alignr(aux2, aux1, offset);
aux1 = aux2;

aux2 = load(input[j-offset+2*packsize]);
_m128 a2 = alignr(aux2, aux1, offset);
aux1 = aux2;

...
\end{code}
\myendfig{alignedloads}{Example of an efficient non-aligned load of sequence of SSE registers.}
\begin{description}
  \item[input] is an array of input data.
  \item[j] represents the first index of the currently processed pack of data.
  \item[offset] is the difference between the actual address of the required data and the preceding aligned address.
\end{description}

Most of our implementation assumes to receive its input by means of more-or-less continuous arrays (e.g., by a sequence of long arrays per every input). Ideally, we would like all these arrays aligned to 16 bytes and be of lengths of multiples of 16 bytes. However, the world is not always ideal. For this reason, we also consider the situation that data inputs are shifted with respect to the output arrays. Since we do not have any means of performing efficient unaligned store\footnote{This is since we would have to store results temporarily somewhere using effects. It is not entirely impossible, but we would have to keep track of these data by means of effects and ensure that the data remaining after termination get written.}, we assume at least ideally aligned output arrays.

As the result, we propose generation of three different instances of the graph-processing code fragment. Use of these code fragments would then be switched during processing in order to achieve efficient processing of all data. These three fragments would be:
\begin{itemize}
\item Non-vectorized realisation.
\item Vectorised realisation utilising unaligned load and store operations.
\item Vectorised realisation utilising aligned load and store operations.
\end{itemize}

The aligned, vectorised realisation would be used on sequences of data rows such that all input and output arrays are continuous across these sequences. Furthermore, these sequences of data would have to be aligned with all output arrays. 

The non-vectorised version would be used to process preambles and epilogues around transitions of data containers. This means processing the last \ttt{offset} rows before an end of any of the containers as well as processing the first \ttt{packsize - offset} after begining of any container. The former processes the last \ttt{offset} lines which cannot be loaded by the vectorised algorithm. The latter ensures that rows processed by vectorised algorithm stay aligned with output containers.

The unaligned vectorised version would be used in cases where no feasible solution using aligned stores exists.





    \section{Buffer marcos} \label{sec:buffer}
\subsubsection{Working principle}

We implement a register-stored SSE buffer by means of macros resolved by the `cfmacros' alias environment. A buffer created by these macros consists of a number of circularly organised registers and two variables which denote the number of stored elements (\ttt{contains}) and the current reading position (\ttt{readat}). We also define some constants which allow us to simply retrieve information such as the total capacity (\ttt{capacity}) of the buffer.

Stores (pushes) are performed by means of a switch which uses the value \ttt{(readat + contains) \% capacity} to determine a code fragment which inserts a single element at the correct position. We assume that the target C compiler optimises such switch into a single jump instruction with its address determined by the expression.  Loads (pops/peeks) are implemented likewise but provide also vectorised version using the \ttt{alignr} instruction. 

Of course, this scheme requires quite large numbers of SSE registers even for small graphs, while standard CPUs provide only 16 such registers. We assume this data to be stored mostly in the L1 cache and automatically loaded when needed. We leave spill code\footnote{The actual machine load and store instructions added after the first iteration of scheduling and register allocation for variables which could not fit into the actual CPU registers.} optimisations to the target C compiler. The compiler may also be supplemented with profiling statistics for this purpose.

\subsubsection{Buffer API}

%API of these buffers follows:

\mybeginfig
\begin{loosecode}
BUFFER_DECL(capacity, vsize, inpacksize, outpacksize, id, type)
BUFFER_SIZE(id)

BUFFER_FULL(id)
BUFFER_FULL_GRANULAR(id)

BUFFER_EMPTY_GRANULAR(id)
BUFFER_EMPTY(id)

BUFFER_PUSH_SIMPLE(capacity, vsize, id, val)
BUFFER_PUSH_ONE(capacity, vsize, id, typeabbrev, val)

BUFFER_CONSUME_ONE(id)
BUFFER_CONSUME_VECTOR(id)

BUFFER_PEEK_SIMPLE(capacity, vsize, id, to)
BUFFER_PEEK_ONE(capacity, vsize, id, typeabbrev, to)
BUFFER_PEEK_VECTOR(capacity, vsize, id, typeabbrev, to)
\end{loosecode}
\myendfig{api10}{Buffer creation macros API reference.}

\begin{description}
  \item \ttt{id} is an identifier of the buffer.
  \item \ttt{capacity} is the total capacity of the buffer.
  \item \ttt{vsize} (vector size) is the number of elements which can be stored into a single SSE register.
  \item \ttt{typeabbrev} is the SSE type suffix.
  \item \ttt{inpacksize and outpacksize} are the sizes of data packs which are processed by a single invocation of the partitions (again the $w$ value of the two partitions the buffer connects). Current implementations assign this number uniformly.
  \item \ttt{DECL} produces declarations of all variables realising the buffer. These are to be accumulated into a layer of a writer which goes into global scope (w.r.t. the code of partitions).
  \item \ttt{SIZE} returns the number of elements actually stored in this container.
  \item \ttt{FULL, EMPTY} resolve to expression which determine whether the buffer is full or empty. The \ttt{GRANULAR} versions take the \ttt{inpacksize} and \ttt{outpacksize} into account, meaning whether an entire vector of values produced by a vectorised partition code can fit into the buffer.
  \item \ttt{PUSH} adds a value into the container.
  \item \ttt{CONSUME} increases the \ttt{readat} counter and decreases the \ttt{contains} counter, effectively removing the head element from the buffer.
  \item \ttt{PEEK} allows reading of the head element (or vector).
  \item \ttt{SIMPLE} versions of the \ttt{PUSH} and \ttt{PEEK} operations do not take the subvector structure into account. A simple variation of the buffer for simple types (such as \ttt{int}) may be constructed using these operations.
\end{description}

\section{ Implementation of splits and merges }

\label{sec:splitsmerges}

As we have already said, the control flow is implemented according to Algorithm \ref{alg:crawler}. For the purpose of generation of all the conditions, we allow operations to have defined an arbitrary number of instruction patterns. Results of these patterns are then accumulated to different writer layers, which are intended for specific chunks of code as shown in the diagram of Algorithm \ref{alg:crawler}. It remains to discuss some details.

Every control flow node gets a unique identifier assigned as its parameter. (This is done during node expansion.) These ids are preserved during node expansion. Every control flow operation uses this id to construct identifiers for all buffers it requires. These buffers are then declared in one of the expanded nodes by the \ttt{DECL} macro and further used by other macros.

Buffer sizes are assigned into parameters with respect to edge position. This means that aliases (node parameters) like \ttt{buffercoefin1}, \ttt{buffercoefin2}, etc. are assigned by the size assignment algorithm (implemented as part of the cf graph transformer in \ttt{cf\_transform.h}) and then used in macro expansions. This may be seen in Figure \ref{fig:balancer}.

Also, every expanded control flow node adds a new label into an accumulator (writer layer) dedicated to labels of the current partition. This way, the expanded nodes maintain links to partitions which contain their data sinks/sources without needing any explicit information about identifiers of these partitions.

Both splits and merges are currently provided only in non-vectorised version -- meaning that the actual split or merge of two streams is performed by a simple loop-like mechanism which always handles one record only. vectorised versions of control flow operations may be easily added once a vectorised version of the \ttt{BUFFER\_PUSH} operation exists.

Despite our brevity, we have described what was worth describing. All other details are technicalities boiling down to the diagram of Algorithm \ref{alg:crawler}. The actual implementation of control flow instructions may be found in \ttt{sse\_set/src\_cf.csv}. Corresponding graph transformations (node expansion, cycle removal and buffer size assignment algorithm) may be found in \ttt{cf\_transform.h}.

\section{ Summary of the provided instruction table }

The C instruction table which we provide may be found in the \ttt{sse\_set} directory of ctb. It covers:
\begin{itemize}
  \item Standard C operators in both vectorised and non-vectorised version.
  \item The discussed data types.
  \item Width conversions.
  \item Load and store operations in both aligned and unaligned versions.
  \item Split and merge operations. These are provided in different variations in directories of the corresponding tests. 
\end{itemize}

The table has been tested by mechanism described in Section \ref{sec:testing}.




    \section{Benchmarks} \label{sec:bench}

\subsubsection{Matrix multiplication}

Scenario (test directory) 16 contains a benchmark of 2-dimensional integer matrix multiplication. This is a graph without any control flow. We used 10000 records. Time was measured by the \ttt{clock()} function. Results are shown in Figure \ref{fig:bench1}. 

The handwritten solution was implemented by nested loops, so its run-time may have been negatively affected by insufficient compiler optimisation. We ran compilation with the \ttt{-O3} flag.

The pack size was 4 since there were no instructions which could use higher vectorisation factor. Thus, load and store instructions were vectorised with factor 4 while multiplications were vectorised with factor 2. (The SSE set does not provide multiplication with higher vectorisation factor than 2.) Also note that the shifted test took significantly longer time due to the cost of its offset switch and an additional preload instruction per every input node. The performance of the shifted version should converge towards efficiency of the aligned solution with increasing pack size (again the $w$ number).

\mybeginfig
\begin{center}
\begin{tabular}{c|c|c|c|c|c}
  \          & handwritten & non-vectorised & aligned & unaligned & shifted \\ \hline
  clang++    & 145          & 53           & 28     & 29       & 46     \\ \hline
  g++        & 124          & 56           & 28     & 30       & 39     \\ 
\end{tabular}
\end{center}
\myendfig{bench1}{Benchmark of two dimensional matrix multiplication (in CPU ticks).}

\subsubsection{Split instruction}

Scenario 17 contains a benchmark of a graph consisting of a single instruction -- a split instruction. This instruction splits a stream of 32-bit integers into two new streams. Since this graph contains bools, the pack size is 16. Other properties remain the same.

We do not provide a vectorised version of the split instruction, so we did not expect good performance. Also, the graph algorithm has to perform two stores and multiple width conversions while the handwritten version just loads and stores the variables without any penalty. Considering this, the results shown in Figure \ref{fig:bench2} are still reasonable ones.

\mybeginfig
\begin{center}
\begin{tabular}{c|c|c|c|c|c}
  \          & handwritten & non-vectorised & aligned & unaligned & shifted \\ \hline
  clang++    & 41          & 150           & 166     & 171       & 166     \\ \hline
  g++        & 36          & 135           & 128     & 152       & 133     \\ 
\end{tabular}
\end{center}
\myendfig{bench2}{Benchmark of the non-vectorised split instruction (in CPU ticks).}


      

%\begin{code}
%for i in abc
%do
%done
%\end{code}
%
%\graph{simple}
%\graph{simple2}
%\graph{cf}
%



%\chapter*{Conclusion}
%\addcontentsline{toc}{chapter}{Conclusion}


%%% Bibliography
%%% Bibliography (literature used as a source)
%%%
%%% We employ bibTeX to construct the bibliography. It processes
%%% citations in the text (e.g., the \cite{...} macro) and looks up
%%% relevant entries in the bibliography.bib file.
%%%
%%% The \bibliographystyle command selects, which style will be used
%%% for references from the text. The argument in curly brackets is
%%% the name of the corresponding style file (*.bst). Both styles
%%% mentioned in this template are included in LaTeX distributions.

%\bibliographystyle{ieeetr}  

%\bibliographystyle{plainnat}    %% Author (year)
\bibliographystyle{unsrt}     %% [number]

\renewcommand{\bibname}{Bibliography}


%%% Generate the bibliography. Beware that if you cited no works,
%%% the empty list will be omitted completely.

\bibliography{bibliography}

%%% If case you prefer to write the bibliography manually (without bibTeX),
%%% you can use the following. Please follow the ISO 690 standard and
%%% citation conventions of your field of research.

% \begin{thebibliography}{99}
%
% \bibitem{lamport94}
%   {\sc Lamport,} Leslie.
%   \emph{\LaTeX: A Document Preparation System}.
%   2nd edition.
%   Massachusetts: Addison Wesley, 1994.
%   ISBN 0-201-52983-1.
%
% \end{thebibliography}


%%% Figures used in the thesis (consider if this is needed)
\listoffigures

%%% Tables used in the thesis (consider if this is needed)
%%% In mathematical theses, it could be better to move the list of tables to the beginning of the thesis.
%\listoftables

%%% Abbreviations used in the thesis, if any, including their explanation
%%% In mathematical theses, it could be better to move the list of abbreviations to the beginning of the thesis.
%\chapwithtoc{List of Abbreviations}

%%% Attachments to the bachelor thesis, if any. Each attachment must be
%%% referred to at least once from the text of the thesis. Attachments
%%% are numbered.
%%%
%%% The printed version should preferably contain attachments, which can be
%%% read (additional tables and charts, supplementary text, examples of
%%% program output, etc.). The electronic version is more suited for attachments
%%% which will likely be used in an electronic form rather than read (program
%%% source code, data files, interactive charts, etc.). Electronic attachments
%%% should be uploaded to SIS and optionally also included in the thesis on a~CD/DVD.
\chapwithtoc{Attachments}
\begin{itemize}
  \item Source code of our framework with examples of instruction tables and example implementations of alias environments should be distributed along with electronic version of this thesis. Plain git repository may be found at \url{http://www.ktweb.cz/ctb}. Browsable mirror is to be found at \url{http://www.github.com/kareltucek/ctb}. Description of directory structure may be found in Section \ref{sec:directories}.
\end{itemize}

\openright
\end{document}
